% =====================================================================
% Paper B: Omitted risk in fitted return models
% Version for submission (VFINAL), 5 November 2026. One self-contained file; references written
% out at the end in alphabetical order. Open items are kept as LaTeX
% comments starting with "% OPEN:" (they do not appear in the PDF).
% =====================================================================
\documentclass[11pt,a4paper]{article}
\usepackage[margin=2.5cm]{geometry}
\usepackage{amsmath,amssymb,amsthm,booktabs,enumitem}
\usepackage[authoryear,round]{natbib}
\usepackage[table]{xcolor}
\usepackage{graphicx}
\usepackage[hidelinks]{hyperref}

% ---------- theorem environments ----------
\newtheorem{theorem}{Theorem}
\newtheorem{lemma}{Lemma}
\newtheorem{proposition}{Proposition}
\newtheorem{corollary}{Corollary}
\theoremstyle{definition}
\newtheorem{assumption}{Assumption}
\theoremstyle{remark}
\newtheorem{remark}{Remark}

% ---------- notation ----------
\emergencystretch=3em
\newcommand{\E}{\mathbb{E}}
\newcommand{\F}{\mathcal{F}}
\newcommand{\tr}{^{\top}}

\title{Omitted risk in fitted return models:\\ displacement, regret, and a robust test for one-sided jumps}
\author{Joilson Giorno\thanks{Graduate Program in Economics, State University of Maring\'a (UEM), Maring\'a, Paran\'a, Brazil. E-mail: \texttt{jgiorno@uem.br}. This paper belongs to a research program on regimes, jumps and heavy tails in return models; the companion paper on regime persistence is \citet{giornoA}.}}
\date{November 2026}

\begin{document}
\maketitle

\begin{abstract}
Return models used for financial risk management are mostly symmetric and continuous, whereas asset returns contain rare and large downside falls. We analyze the econometric and economic consequences for a fitted model when such one-sided jumps are omitted. For any parametric model, correctly specified or not, an omitted feature of size $\varepsilon$ shifts the pseudo-true parameters by $\varepsilon H_M^{-1}\E[S_MW]$, where $W$, the score of the omitted feature, is the same for every fitted model. For omitted one-sided jumps, $W$ has a single-impulse form that is easy to evaluate by simulation; the resulting error in value-at-risk (VaR) is of first order in $\varepsilon$, while the quantile-loss regret of an otherwise correct model is of second order, with an explicit coefficient. Omitted downside risk is therefore almost harmless for mean forecasts but costly for tail risk: across 4{,}000 simulated paths, jumps on 2\% of the days leave 99\% VaR 5--7\% too low and raise the frequency of Basel yellow-zone outcomes from 11--14\% to 21--22\%. We propose a test for one-sided impulses that remains valid when the fitted model has the wrong tails, scale or volatility dynamics. The test function is odd in the standardized residual and, after a Jacobian correction, Neyman-orthogonal to all estimated parameters. It needs no long-run variance estimator, and its local power is governed by the same single-impulse functional that drives the parameter displacement. In simulations the test keeps its nominal size (4.3--4.5\%) and is more powerful than skewness and omnibus symmetry tests, whereas a naive score test rejects up to 99.6\% of jump-free paths. Fisher information bounds show that moderate jumps can be detected but not estimated, which leads to a sequential procedure for risk managers. On long daily samples of SPY, the S\&P 500 and the Ibovespa, the test detects downside asymmetry in all three series; this asymmetry behaves like a skewed innovation law rather than rare jumps, and modeling it restores the out-of-sample calibration of 99\% VaR for the two U.S.\ series.
\end{abstract}

\noindent\textbf{Keywords:} jumps; misspecification; local sensitivity; score test; conditional symmetry; hidden Markov models; value-at-risk.\\
\textbf{JEL:} C12, C22, C52, C58, G17.

\section{Introduction}\label{sec:intro}

Banks, asset managers and financial regulators measure market risk by fitting econometric models to return series. The models most often used assume symmetric noise structures, such as GARCH-type processes with Student-$t$ innovations \citep{bollerslev1986,ding1993,gjr1993}, sometimes augmented with Markov regime-switching dynamics \citep{hamilton1989,haas2004}. Their risk forecasts have direct economic use in allocation decisions. More importantly, under the Basel rules, the market-risk capital of a bank is a multiple of its 99\% value-at-risk (VaR), and the multiple increases when the backtest finds too many exceptions \citep{basel1996}. However, financial returns are prone to sudden and severe downside shocks, or jumps, driven by market crashes, currency devaluations or sudden liquidity contractions \citep{merton1976,eraker2003}. Standard continuous models are not designed to accommodate these one-sided discontinuities. Understanding how fitted models absorb this omitted downside risk, evaluating its economic cost, and detecting its presence are the central objectives of this paper.

When symmetric econometric models are fitted to jump-contaminated return series, the extreme downside observations do not simply vanish. Because the maximum-likelihood estimator must explain every data point, the omitted downward jumps are absorbed by the parameters capable of imitating them: the estimated location shifts downward, the scale parameters inflate and the estimated degrees of freedom fall towards heavier tails. The fitted specification may look plausible, yet it describes a different risk from the true one. This paper addresses three questions about this misspecification: (i) where does the omitted downside risk go in the parameter space? (ii) what is the cost to the user of the model? and (iii) can risk managers detect omitted jump risk without knowing the true data-generating process?

The answer to the second question gives the main economic message of this paper: the cost of the omission depends on the loss function of the decision-maker. For mean forecasting, omitting rare jumps generates a negligible loss. For instance, in a series with unit variance, downward jumps on 2\% of the trading days with a mean size of 1.8 produce a squared-error forecast regret of only $0.0006$. For tail-risk management, however, the same jump process is costly. In our Monte Carlo evaluations, ignoring these jumps leaves 99\% VaR forecasts 5--7\% too low, raises the probability of entering the Basel yellow zone from 11--14\% to 21--22\%, and leaves regulatory capital 4--5\% below what the true risk requires. Consequently, a model can perform well for point forecasting while failing in tail-risk measurement. The framework developed here, which combines a local displacement theory with a robust specification test, gives regulators and risk managers practical tools to address this weakness.

\paragraph{Contributions.} This paper makes four contributions to the econometrics of financial risk.

\emph{First, a local theory of parameter displacement under misspecification.} For any parametric model, correctly specified or not, an omitted feature of size $\varepsilon$ shifts the pseudo-true parameter vector by $\varepsilon H_M^{-1}\E[S_MW]$ (Theorem~\ref{thm:displacement}, Proposition~\ref{prop:general}). Here $W$ is the score of the omitted feature under the true data-generating process. Crucially, $W$ does not depend on the fitted model: all candidate models are pushed along one common direction, and each responds through its own score $S_M$ and Hessian $H_M$ (Corollary~\ref{cor:common}). For omitted one-sided jumps, $\E[S_MW]$ is the expected change in the model's score when one single impulse is added to the path, which is easy to evaluate by simulation. The result holds for a fixed sample and, under summability conditions on the effect of an impulse, for the probability limit of the quasi-maximum-likelihood estimator (Theorem~\ref{thm:largeT}).

\emph{Second, the economic cost of omitted risk.} The displacement expansion gives the first-order error in model-implied VaR (Corollary~\ref{cor:var}). Combined with the pinball (quantile) loss, the loss function consistent with VaR \citep{koenker1978,gneiting2011}, it gives the economic regret of the omission. For a model that is correctly specified without jumps, the regret is of second order in the jump intensity, $\tfrac12\varepsilon^2\E[f_t\dot\Delta_t^2]$, with an explicit coefficient (Corollary~\ref{cor:regret}). This explains why omitted jumps are almost harmless for mean forecasts but costly for tail risk (Remark~\ref{rem:meantail}). Monte Carlo experiments confirm the predicted orders: doubling the jump intensity increases the VaR shortfall by a factor of 1.85--1.93 and the pinball-loss regret by a factor of about 3.4, close to the theoretical values of 2 and 4.

\emph{Third, a score test robust to symmetric misspecification of tails and dynamics.} We formulate a score test of zero jump intensity within the class of Bernoulli--Gamma jumps ($-B_tG_t$), made robust by taking the difference between the downward- and upward-impulse scores (Section~\ref{sec:test}). Because this difference is odd in the standardized residual, it has conditional mean zero under any conditionally symmetric law, whether or not the fitted model has the right tails, scale or volatility dynamics (Theorem~\ref{thm:symm}). A Jacobian projection makes the test Neyman-orthogonal to all estimated parameters (Lemma~\ref{lem:mds}). Because the test terms form a martingale difference sequence, no long-run variance (HAC) estimator is required, and critical values over the unidentified jump scale are obtained from a Gaussian supremum or a multiplier bootstrap (Theorem~\ref{thm:null}). The local power of the test is governed by the same single-impulse functional that determines the parameter displacement (Theorem~\ref{thm:local}).

\emph{Fourth, a decision boundary between detecting jumps and estimating them.} Fisher information bounds for the jump parameters (Proposition~\ref{prop:jumpinfo}) show two regimes. Large jumps provide information in proportion to the number of jump events. Moderate jumps overlap with the heavy tails of the innovation distribution, and the information about their scale and shape is only of order $\varepsilon^2$, the same order as the noncentrality of the test. Consequently, moderate downside jumps can be detected by the test but cannot be estimated with acceptable precision. This leads to a sequential procedure for market risk management (Section~\ref{sec:procedure}): test first, check whether the sample contains enough information to estimate the jumps, and if not, report the displacement bias of the symmetric model instead of fitting an unstable jump model.

One functional links the two halves of the paper: the single-impulse functional that determines how far an omitted jump moves the parameter estimates also determines how visible the jump is to the test. Our simulations show the economic implication of this link. In data-generating processes whose regimes differ strongly in variance, symmetric models absorb omitted jumps most completely into their volatility and tail parameters. The VaR shortfall is largest in these settings (11.9\% for GARCH-$t$ in one design), and the jumps are hardest to detect (power 40--63\%). Omitted tail risk is most costly precisely where it is hardest to detect.

\paragraph{Empirical findings.} We apply the sequential procedure to long daily return series of the SPY ETF (1993--2026), the S\&P 500 index (1990--2026) and the Brazilian Ibovespa index (1995--2026, after the Real Plan). The robust test rejects conditional symmetry towards downward impulses in all three series, and towards upward impulses in none. However, the downside asymmetry does not behave like rare and large jumps: a fitted jump model turns it into frequent small impulses, which is in effect a skewed innovation distribution, and the information pre-check advises against reporting jump parameters. Modeling the downside asymmetry, with either a skew-$t$ innovation law or a jump component, restores the out-of-sample calibration of 99\% VaR for the two U.S.\ series over 2020--2026. For the Ibovespa, whose symmetric VaR was already conservative over that period, modeling the asymmetry does not help out of sample. For risk managers and supervisors, the relevant practical question is therefore not ``are there discrete jumps?'' but ``is there unmodeled downside asymmetry?'', and our procedure answers it without requiring knowledge of the true specification.

\subsection{Related literature}\label{sec:lit}

\paragraph{Misspecification and local sensitivity.} The analysis of misspecified models goes back to \citet{white1982}, who showed that quasi-maximum-likelihood estimators converge to pseudo-true parameter values. A growing literature studies parameter sensitivity under small perturbations of a model. \citet{kitamura2013} study estimation in infinitesimal neighborhoods of moment conditions, \citet{andrews2017} measure the sensitivity of parameter estimates to the moments used, \citet{chen2024} develop inference for moment models whose misspecification is bounded by a statistical divergence, and \citet{bonhomme2022} construct estimators that minimize the worst-case bias within a neighborhood of the model. Bonhomme and Weidner also link the size of the neighborhood to the power of a specification test, close to the idea of our Remark~\ref{rem:onefunctional}. Our approach is complementary: instead of modifying the estimator, we take the models used in financial practice and determine where one specific, economically important omission drives their parameter estimates. Because the omitted feature is specified, the displacement has a closed single-impulse form, and its cost can be measured directly in the loss function relevant for the decision. This agrees with the view of \citet{andrews2026} that pseudo-true parameters matter because they feed decisions; in our setting, the decision is the capital charge implied by VaR.

\paragraph{Jumps in financial returns.} Jump models in finance start with \citet{merton1976}; \citet{kou2002} introduced double-exponential jumps with separate upward and downward components, estimated by maximum likelihood in \citet{ramezani2007}. Discrete-time GARCH models with jumps were developed by \citet{chan2002}, \citet{maheu2004} and \citet{duan2006}. In 72 years of daily Dow Jones returns, \citet{chan2002} find that the conditional jump intensity varies strongly over time and that the mean jump size depends on the sign of the previous return. Evidence of jumps in prices and volatility also comes from continuous-time models \citep{eraker2003} and from nonparametric tests on intraday data \citep{bns2006,leemykland2008}. Our test serves a different purpose: it uses daily or lower-frequency returns, runs within any fitted symmetric model, and does not assume that this model is correct. Its target is also specific: because the test function is the downward-impulse score minus the upward-impulse score, it measures the net downward jump component, and equal upward and downward jumps cancel at first order (Remark~\ref{rem:twosided}). This is the component that matters for downside tail risk.

\paragraph{Regimes and specification tests in Markov-switching models.} \citet{hamilton1996} showed that the score of a Markov-switching model can be computed directly from the smoothed regime probabilities, and used this to build Lagrange multiplier and Newey--Tauchen--White tests for omitted autocorrelation, omitted ARCH effects, misspecified Markov dynamics and omitted explanatory variables. The jump direction $W_T$ in \eqref{eq:W} has the same structure: it weights the jump-incidence ratio of each state by the smoothed probability of that state. However, Hamilton's tests require the Markov-switching model to be correctly specified under the null, while our test remains valid when the fitted model has the wrong tails, scale or dynamics. The asymptotic theory of the maximum-likelihood estimator in these models is given by \citet{douc2004}. Testing when a nuisance parameter is not identified under the null, here the jump scale, follows \citet{davies1977,davies1987}, \citet{andrews1994ap}, \citet{hansen1992} and \citet{hansen1996}. The companion paper \citet{giornoA} studies when regime persistence can be estimated in Student-$t$ hidden Markov models; we use its effective-sample criterion to choose the base model in the first step of our procedure.

\paragraph{Symmetry, skewness and asymmetric distributions.} Tests of conditional symmetry include \citet{bainh2001}, \citet{su2006}, \citet{fernandes2015} and \citet{chen2024cs}. \citet{bainh2005} propose tests for skewness and kurtosis in time series, and \citet{maasoumi2009} develop an entropy-based test that compares the kernel density of a series with the density of its reflection about the mean, with bootstrap critical values for dependent data. Their statistic, $S_\rho=\tfrac12\int(f_1^{1/2}-f_2^{1/2})^2$, is the Hellinger-type metric entropy; it equals one half of the Tsallis relative entropy of order $1/2$ between the two densities, a link with the Tsallis family of heavy-tailed laws to which our Student-$t$ specifications belong. These symmetry tests are omnibus: they look for asymmetry anywhere in the distribution. \citet{maasoumi2009} note that moment-based tests gain power when they use the one-sided direction of the asymmetry; our test is built on that direction from the start. Moreover, when omnibus tests are applied to model residuals, the estimated location changes their null distribution (Section~\ref{sec:b12}). Our test directs its power at downside impulses, has a closed-form local power, and is orthogonal to the estimated parameters without a martingale transformation. In our simulations it is more powerful than a skewness test and an omnibus symmetry test, both against jumps and against skewed innovation laws (Section~\ref{sec:b12}). Skewed innovation distributions for GARCH models \citep{hansen1994,fernandez1998,trottier2016} are the natural alternative explanation of a rejection; we compare them with jump models by the non-nested likelihood-ratio test of \citet{vuong1989}.

\paragraph{Evaluation of value-at-risk.} VaR forecasts are evaluated by coverage backtests \citep{kupiec1995,christoffersen1998,christoffersen2004}, by loss functions related to the regulatory rules \citep{lopez1999}, by quantile regression \citep{gaglianone2011} and by consistent scoring functions \citep{gneiting2011}. \citet{kuester2006} compare many VaR models out of sample. \citet{escanciano2010} show that standard backtests ignore estimation risk, and model risk when the model is misspecified, so that their critical values can be wrong; the distortion grows with the ratio of the evaluation sample to the estimation sample. These studies evaluate the accuracy of a given forecast. Our framework predicts the first-order size of the VaR error caused by an omitted feature, and how the tail loss scales with the jump intensity.

\paragraph{Outline.} Section~\ref{sec:setting} describes the setting. Section~\ref{sec:where} derives the displacement theory and the economic regret. Section~\ref{sec:test} constructs the robust test and gives its asymptotic properties. Section~\ref{sec:after} studies when the jumps can be estimated and states the sequential procedure. Section~\ref{sec:mc} presents the Monte Carlo evidence, Section~\ref{sec:data} the application to equity returns, and Section~\ref{sec:conclusion} concludes. Proofs of Theorems~\ref{thm:largeT} and~\ref{thm:local} and of Proposition~\ref{prop:jumpinfo} are in Appendix~\ref{app:proofs}, and additional simulation evidence is in Appendix~\ref{app:extra}.

\section{Setting}\label{sec:setting}
\paragraph{Data-generating process.} Let $X=(X_1,\dots,X_T)$ be a strictly stationary return process. In the simulations, $X$ is a centered two-state Student-$t$ hidden Markov model with parameters $\psi_0=(p_{00},p_{11},\nu_0,\nu_1,s_0,s_1)$: given the state $K_t=j$, $X_t$ has the Student-$t$ density $f_j$ with $\nu_j$ degrees of freedom and scale $s_j$, and $K_t$ is a two-state Markov chain. The observed returns are contaminated by rare downward impulses,
\begin{equation}\label{eq:dgp}
Y_t=X_t-B_tG_t,\qquad B_t\sim\text{Bernoulli}(\varepsilon),\quad G_t\sim\text{Gamma}(k,\theta),
\end{equation}
with $B$, $G$ and $X$ mutually independent; $\theta$ is the Gamma scale and $k\theta$ the mean jump size. We call these \emph{Bernoulli--Gamma jumps} ($-B_tG_t$): $B_t$ decides whether a jump occurs, $G_t$ how large it is, and the jump is always downward. Write $h_j(y)=\E f_j(y+G)$ for the state-$j$ density of a jump-contaminated observation and $p_\varepsilon(y_{1:T})$ for the joint density of $Y$.

\paragraph{The jump direction.} Let $\gamma_{tj}(y)=\Pr_0(K_t=j\mid Y_{1:T}=y)$ be the jump-free smoothing probabilities. The jump-incidence score is
\begin{equation}\label{eq:W}
W_T(y)=\partial_{\varepsilon+}\log p_\varepsilon(y)\big|_{\varepsilon=0}=\sum_{t=1}^T\sum_{j=0}^1\gamma_{tj}(y)\left\{\frac{h_j(y_t)}{f_j(y_t)}-1\right\}.
\end{equation}
$W_T$ depends only on the data-generating process, not on any fitted model (Appendix~\ref{app:steps}). As in \citet{hamilton1996}, this score is a byproduct of the smoothed probabilities.

\paragraph{Fitted models.} A forecaster fits a parametric model $M$ with log-likelihood $\ell_M(\beta;y_{1:T})$, $\beta\in B\subset\mathbb R^m$, without a jump component. $M$ may be misspecified even when $\varepsilon=0$: a static Student-$t$, GARCH-$t$ \citep{bollerslev1986}, APARCH-$t$ \citep{ding1993}, or the hidden Markov model itself. Write $\beta^*_{M,T}(\varepsilon)$ for a maximizer of $Q_{M,T}(\beta,\varepsilon)=\E_\varepsilon\ell_M(\beta;Y)$, $S_M=\nabla_\beta\ell_M(\beta^*_{M,T}(0);Y)$ and $H_{M,T}=-\E_0\nabla^2_\beta\ell_M(\beta^*_{M,T}(0);Y)$.

\section{Where omitted risk goes}\label{sec:where}
\begin{theorem}[Omitted-jump displacement in any fitted model]\label{thm:displacement}
Fix $T$, $k$ and $\theta$. Suppose $\beta^*_{M,T}(0)$ is interior and is the locally unique maximizer of $Q_{M,T}(\cdot,0)$; $\ell_M$ is twice continuously differentiable in $\beta$ with integrable dominators for $\nabla\ell_M$ and $\nabla^2\ell_M$ under $p_\varepsilon$ for $\varepsilon$ in a right neighborhood of $0$, and $\E_0\|S_MW_T\|<\infty$; and $H_{M,T}$ is positive definite. Then there is a locally unique branch of maximizers with
\begin{equation}\label{eq:displacement}
\begin{aligned}
\beta^*_{M,T}(\varepsilon)&=\beta^*_{M,T}(0)+\varepsilon H_{M,T}^{-1}B_{M,T}+o(\varepsilon),\\
B_{M,T}&=\E_0[S_MW_T]=\sum_{t=1}^T\E_0\E_G\big\{S_M(X^{[t,G]})-S_M(X)\big\},
\end{aligned}
\end{equation}
where $X^{[t,g]}$ replaces $X_t$ by $X_t-g$ and $S_M(\cdot)$ is evaluated at $\beta^*_{M,T}(0)$.
\end{theorem}
\begin{proof}
Let $F(\beta,\varepsilon)=\E_\varepsilon\nabla_\beta\ell_M(\beta;Y)=\int\nabla_\beta\ell_M(\beta;y)\,p_\varepsilon(y)\,dy$. Because $p_\varepsilon$ is a polynomial in $\varepsilon$ whose coefficients are densities or differences of densities (Appendix~\ref{app:steps}, step 1), the dominators allow differentiation under the integral, so $F$ is continuously differentiable on $B\times[0,\bar\varepsilon)$ and extends to a polynomial in $\varepsilon$ for $\varepsilon$ slightly negative. At $\varepsilon=0$, $F(\beta^*_{M,T}(0),0)=0$ by the first-order condition. Its derivatives are $D_\beta F=-H_{M,T}$ and $\partial_\varepsilon F=\int S_M\dot p_0\,dy=\E_0[S_MW_T]$, since $\dot p_0=p_0W_T$ (step 2). The second representation of $B_{M,T}$ follows from step 3. Because $H_{M,T}$ is nonsingular, the implicit-function theorem gives a unique continuously differentiable branch with $0=B_{M,T}-H_{M,T}\dot\beta$, i.e.\ $\dot\beta=H_{M,T}^{-1}B_{M,T}$. No information identity is used, so $M$ need not be correctly specified at $\varepsilon=0$. When $M$ is the HMM itself, $H_{M,T}$ is its Fisher information.
\end{proof}

\begin{proposition}[Any smooth omitted feature]\label{prop:general}
Let $\{p_\varepsilon\}$ be any family of laws for $Y_{1:T}$ with $p_0$ the law without the omitted feature, differentiable at $\varepsilon=0$ with $\dot p_0=p_0W$ for a statistic $W$ with $\E_0W=0$, and such that the dominators of Theorem~\ref{thm:displacement} hold for $\varepsilon$ in a neighborhood of zero. Then, under the remaining conditions of Theorem~\ref{thm:displacement},
\[
\beta^*_{M,T}(\varepsilon)=\beta^*_{M,T}(0)+\varepsilon H_{M,T}^{-1}\E_0[S_MW]+o(\varepsilon).
\]
\end{proposition}
\begin{proof}
The proof of Theorem~\ref{thm:displacement} uses the jump structure only to establish differentiability of $\varepsilon\mapsto p_\varepsilon$ (step~1) and the form of $\dot p_0$ (step~2); both are now assumed.
\end{proof}
\noindent The coefficients are specific to the omitted feature: only $W$ changes. Examples: an omitted drift $p_\varepsilon(y)=p_0(y-\varepsilon m)$ gives $W=-m\tr\nabla_y\log p_0(y)$; omitted one-sided jumps give $W=W_T$ in \eqref{eq:W}; an omitted skewing of the innovation density gives the score of the skewing parameter. The result is local: it describes small omissions, and its accuracy for a given $\varepsilon$ must be checked (Table~\ref{tab:cross}).

\begin{corollary}[Common jump direction]\label{cor:common}
$W_T$ is the same for every fitted model. The omitted jump therefore acts as one fixed direction in the space of square-integrable statistics, and each model's displacement is the $H_M^{-1}$-weighted image of the projection of that direction on the model's score space. Models whose scores are uncorrelated with $W_T$ do not move at first order, whatever their fit.
\end{corollary}

\begin{corollary}[Local VaR displacement]\label{cor:var}
For any differentiable functional $\phi(\beta)$, such as the model-implied lower quantile $Q_M(\alpha;\beta)$, $\partial_\varepsilon\phi(\beta^*_{M,T}(\varepsilon))|_{0+}=\nabla\phi\tr H_{M,T}^{-1}B_{M,T}$. For the stationary marginal quantile of the data-generating process, $\partial_\varepsilon Q^{\rm true}_\alpha|_{0+}=\{F_0(Q_\alpha)-\E_GF_0(Q_\alpha+G)\}/f_0(Q_\alpha)$, exactly for any $\theta$. The difference of the two is the first-order VaR shortfall of model $M$ per unit of jump intensity.
\end{corollary}
\begin{proof}
The first statement is the chain rule applied to \eqref{eq:displacement}. For the second, $\Pr(Y_t\le q)=(1-\varepsilon)F_0(q)+\varepsilon\,\E_GF_0(q+G)$; differentiate $\Pr(Y_t\le Q_\alpha(\varepsilon))=\alpha$ at $\varepsilon=0$.
\end{proof}

\begin{corollary}[Value-at-risk regret]\label{cor:regret}
Let $\rho_\alpha(e)=(\alpha-\mathbf 1\{e<0\})e$ be the quantile loss, $Q_t(\varepsilon)$ the true conditional $\alpha$-quantile of $Y_t$ given $\F_{t-1}$, with conditional density $f_t$ positive and continuous at $Q_t(0)$, and $q_{M,t}(\varepsilon)=q_{M,t}(\beta^*_{M,T}(\varepsilon))$ the model's forecast. Then
\[
\E_\varepsilon\big[\rho_\alpha(Y_t-q_{M,t})-\rho_\alpha(Y_t-Q_t)\big]=\E_\varepsilon\int_{Q_t}^{q_{M,t}}\big\{F_t(x)-\alpha\big\}\,dx\;\ge0 .
\]
If the model is correctly specified at $\varepsilon=0$, so that $q_{M,t}(0)=Q_t(0)$, and $\Delta_t(\varepsilon)=q_{M,t}(\varepsilon)-Q_t(\varepsilon)$ is differentiable at $0$ with $\dot\Delta_t=\nabla_\beta q_{M,t}\tr H_{M,T}^{-1}B_{M,T}-\dot Q_t$, then
\[
\text{Regret}_t(\varepsilon)=\tfrac12\,\varepsilon^2\,\E_0\big[f_t(Q_t)\,\dot\Delta_t^2\big]+o(\varepsilon^2).
\]
\end{corollary}
\begin{proof}
The identity follows from $\partial_q\E[\rho_\alpha(Y_t-q)\mid\F_{t-1}]=F_t(q)-\alpha$, which vanishes at $Q_t$; the integrand is nonnegative on the interval of integration. When $\Delta_t(0)=0$, the interval has length $\varepsilon|\dot\Delta_t|+o(\varepsilon)$ and on it $F_t(x)-\alpha=f_t(Q_t)(x-Q_t)+o(\varepsilon)$, which gives $\tfrac12f_t(Q_t)\varepsilon^2\dot\Delta_t^2$; the change of measure from $\E_0$ to $\E_\varepsilon$ affects only higher-order terms.
\end{proof}
\noindent Coverage errors are of first order in $\varepsilon$ (Corollary~\ref{cor:var}), whereas the quantile-loss regret of a model that is correct without the omission is of second order. In the simulation designs this covers the HMM-6 fitted to HMM data. For models already misspecified at $\varepsilon=0$, the baseline regret is positive and the omission changes it at first order; Section~\ref{sec:var} illustrates both cases.
% OPEN (A7): first-order formula for the incremental regret of misspecified models.

\begin{remark}[Point forecasting versus tail-risk management]\label{rem:meantail}
Under squared loss, the regret of a forecast $a_t$ relative to the true conditional mean $m_t$ is $\tfrac12\E[(a_t-m_t)^2]$, so a zero-mean forecast has regret $\tfrac12\E[m_t^2]$. Independent jumps shift the conditional mean of a centered process by $-\varepsilon k\theta$, so the regret of ignoring them in a mean forecast is $\tfrac12(\varepsilon k\theta)^2$: about $0.0006$ for jumps on 2\% of the days ($\varepsilon=0.02$) with mean size 1.8, in data whose daily return variance is of order one. The same jumps leave 99\% VaR 5--7\% too low (Table~\ref{tab:var}). A model judged by its mean-squared forecast error alone therefore gives no warning about its tail risk.
\end{remark}

\subsection{Large samples}\label{sec:largeT}
Theorem~\ref{thm:displacement} holds for a fixed window. What a forecaster obtains is the probability limit of the estimator, which we now characterize. Write the fitted log-likelihood through its one-step contributions, $\ell_M(\beta;y_{1:T})=\sum_t\ell_t(\beta)$, where $\ell_t(\beta)=\ell(\beta;y_t,y_{t-1},\dots)$ is the stationary one-step log predictive density of model $M$ and $s_t=\nabla_\beta\ell_t$. For a finite set (or a random set) $J$ of dates before or at $t$, let $X^{[J,G]}$ denote $X$ with $X_u$ replaced by $X_u-G_u$ for $u\in J$, and define the impulse effects
\[
\Delta_t(J)=s_t\big(X^{[J,G]}\big)-s_t(X).
\]
Let $q(\beta,\varepsilon)=\E_\varepsilon\,\ell_t(\beta;Y)$, the limit of $T^{-1}\E_\varepsilon\ell_M(\beta;Y_{1:T})$.

\begin{assumption}[Large-sample displacement]\label{A:largeT}
(i) $X$ is strictly stationary and ergodic, and the jump variables are i.i.d.\ and independent of $X$, so $Y$ is stationary and ergodic for each $\varepsilon$. (ii) $B$ is compact, $q(\cdot,0)$ has a unique maximizer $\bar\beta^*_M$ in the interior, $\ell_t$ is twice continuously differentiable on a neighborhood $N$ of $\bar\beta^*_M$, and $\E_\varepsilon\sup_B|\ell_t|$, $\E_\varepsilon\sup_N\|s_t\|$ and $\E_\varepsilon\sup_N\|\nabla_\beta s_t\|$ are finite for $\varepsilon\in[0,\bar\varepsilon]$; $\bar H_M=-\E_0\nabla_\beta s_t(\bar\beta^*_M)$ is positive definite. (iii) Impulse effects are summable: there are nonnegative sequences $(a_j)$ and $(c_j)$ with $\sum_ja_j<\infty$ and $\sum_jc_j<\infty$ such that, for every set $J$ of past dates,
\[
\E\sup_N\big|\ell_t(X^{[J,G]})-\ell_t(X)\big|+\E\sup_N\big\|\nabla_\beta s_t(X^{[J,G]})-\nabla_\beta s_t(X)\big\|\le\sum_{j:\,t-j\in J}a_j,
\]
\[
\E\sup_N\Big\|\Delta_t(J)-\sum_{u\in J}\Delta_t(\{u\})\Big\|\le\sum_{\substack{u<v\\u,v\in J}}c_{t-u}\,c_{t-v},\qquad \E\sup_N\|\Delta_t(\{t-j\})\|\le a_j .
\]
\end{assumption}
\noindent Condition (iii) says that the effect of a jump on later contributions fades, and that two jumps interact only weakly. For a static model only $j=0$ matters. For GARCH- and APARCH-$t$ models the conditional scale forgets past residuals geometrically, and the Student-$t$ scores are bounded in the standardized residual, which gives geometric $a_j$ and $c_j$ on a neighborhood where the persistence parameter is below one. For the HMM, the predictive filter forgets its past exponentially under positive transition probabilities and emission densities bounded away from zero on compacts \citep{legland2000,douc2004}, which gives the same geometric decay.

\begin{theorem}[Large-sample displacement]\label{thm:largeT}
Under Assumption~\ref{A:largeT}, for all sufficiently small $\varepsilon>0$, $q(\cdot,\varepsilon)$ has a unique maximizer $\bar\beta^*_M(\varepsilon)$ in $N$, and
\[
\begin{aligned}
\bar\beta^*_M(\varepsilon)&=\bar\beta^*_M+\varepsilon\,\bar H_M^{-1}\bar b_M+o(\varepsilon),\\
\bar b_M&=\sum_{j\ge0}\E_0\E_G\big[\Delta_t(\{t-j\})\big]=\E_0\E_G\Big[\sum_{u\ge t}\big\{s_u(X^{[t,G]})-s_u(X)\big\}\Big],
\end{aligned}
\]
with the scores evaluated at $\bar\beta^*_M$. If in addition $T^{-1}\sum_t\ell_t(\beta)$ computed from the observed sample converges to $q(\beta,\varepsilon)$ uniformly on $B$ almost surely, the quasi-maximum-likelihood estimator satisfies $\hat\beta_T\to\bar\beta^*_M(\varepsilon)$ almost surely.
\end{theorem}
\noindent The proof, in Appendix~\ref{app:proofs}, expands the expected score in $\varepsilon$ using the summability of impulse effects, shows that the maximizer moves continuously, and linearizes the first-order condition.

\begin{remark}[Link with Theorem~\ref{thm:displacement} and computation]\label{rem:largeT}
Under the same summability, $T^{-1}H_{M,T}\to\bar H_M$ and $T^{-1}B_{M,T}\to\bar b_M$, since $T^{-1}B_{M,T}$ is a Ces\`aro average of the impulse sums in $\bar b_M$ with edge effects of order $T^{-1}\sum_jj\,a_j$ when that sum is finite; the finite-window and large-sample displacements then agree to first order. The single-impulse form makes $\bar b_M$ cheap to compute: place an impulse at a random date of a long simulated path and record the change in the model's score over the following window. This is how the predictions in Table~\ref{tab:cross} were computed, with 20{,}000 impulses on a 400{,}000-observation path, instead of fitting each model to thousands of simulated jump paths.
\end{remark}

\begin{lemma}[A jump law cannot be symmetric]\label{lem:noneq}
Let the marginal law of $X_t-m$ be a scale mixture of centered normals (this covers Student-$t$ laws and every mixture of centered Student-$t$ laws, hence every centered Student-$t$ HMM), and let $Y_t=X_t-B_tG_t$ with $\varepsilon>0$ and $G_t\ge0$ nondegenerate with unbounded support. Then the marginal law of $Y_t$ is not symmetric about any finite point. In particular, no location shift of a symmetric model can reproduce it.
\end{lemma}
\begin{proof}
The characteristic function $\varphi_X$ of $X_t-m$ is real, even and strictly positive. If $Y_t$ were symmetric about $c$, then $e^{-iu(c-m)}\varphi_X(u)\varphi_J(u)$ would be real and even, with $J=-B_tG_t$; dividing by $\varphi_X>0$, $J-(c-m)$ would be symmetric about zero, which is impossible for a law bounded above and unbounded below.
\end{proof}
\noindent Lemma~\ref{lem:noneq} explains why the displacement in Theorem~\ref{thm:displacement} is not confined to the location: a symmetric model cannot reproduce the jump by shifting, so it also uses scale and tail parameters.

\section{A test robust to tail misspecification}\label{sec:test}
\subsection{Construction}
Let $\F_t=\sigma(y_s,s\le t)$. The fitted model has symmetric innovations,
\begin{equation}\label{eq:model}
f_t(y;\beta)=\frac{1}{\sigma_t(\beta)}\,g\!\left(\frac{y-\mu_t(\beta)}{\sigma_t(\beta)};\eta\right),\qquad u_t(\beta)=\frac{y_t-\mu_t(\beta)}{\sigma_t(\beta)},
\end{equation}
where $\mu_t$ and $\sigma_t>0$ are $\F_{t-1}$-measurable, $g(\cdot;\eta)$ is a positive density symmetric about zero, and $\eta$, the shape parameter of $g$ that controls the thickness of its tails, is part of $\beta$.\footnote{For the Student-$t$ density, $\eta=\nu$, the degrees of freedom. In the Tsallis parametrization the Student-$t$ law is the $q$-Gaussian with entropic index $q=(\nu+3)/(\nu+1)$, so a smaller $\nu$ corresponds to a larger $q$ and heavier tails.} Static Student-$t$, GARCH-$t$, GJR-$t$ \citep{gjr1993} and APARCH-$t$ with constant mean are of this form, and so is the one-step predictive law of a centered Student-$t$ HMM. Write $\psi(u;\eta)=-\partial_u\log g(u;\eta)$ (odd).

Let $Z\ge0$ have a fixed density (Gamma$(2,1)$ in the implementation), $G=\theta Z$, and $\Theta=[\theta_L,\theta_U]\subset(0,\infty)$. Define
\[
h^-_t(y)=\E f_t(y+G),\qquad h^+_t(y)=\E f_t(y-G),\qquad a_t(\theta)=\frac{h^-_t(y_t)-h^+_t(y_t)}{f_t(y_t)} .
\]
$h^-_t$ is the density of the fitted law with a downward impulse added, $h^+_t$ with an upward one, so $a_t$ is the difference of the two jump-incidence scores. In standardized units, with $c_t=\theta/\sigma_t$,
\begin{equation}\label{eq:a}
a_t(\theta;\beta)=A\big(u_t(\beta);c_t,\eta\big),\qquad A(u;c,\eta)=\frac{\E\,g(u+cZ;\eta)-\E\,g(u-cZ;\eta)}{g(u;\eta)},
\end{equation}
and $A$ is odd in $u$ for every $(c,\eta)$.

\begin{theorem}[Robustness to symmetric misspecification]\label{thm:symm}
Suppose the true conditional law of $y_t$ given $\F_{t-1}$ is symmetric about $\mu_t$, with any shape, and $f_t$ is symmetric about the same $\mu_t$. Then $\E[a_t(\theta)\mid\F_{t-1}]=0$ for every $\theta$ for which the expectation exists, even if $f_t$ has the wrong tails or the wrong scale.
\end{theorem}
\begin{proof}
Symmetry of $f_t$ gives $h^-_t(\mu_t+u)=\E f_t(\mu_t+u+G)=\E f_t(\mu_t-u-G)=h^+_t(\mu_t-u)$, and $f_t(\mu_t+u)=f_t(\mu_t-u)$, so $a_t(\mu_t+u)=-a_t(\mu_t-u)$: $a_t$ is odd about $\mu_t$, and its expectation under any law symmetric about $\mu_t$ is zero.
\end{proof}
\noindent The naive one-sided score $h^-_t/f_t-1$ has mean zero only if $f_t$ is the true law, so it rejects when the tails are merely misspecified (Section~\ref{sec:mc}).

\paragraph{Orthogonalization.} Define the location-direction score $s^\mu_t(\beta)=\sigma_t^{-1}\psi(u_t)\nabla_\beta\mu_t$ and, at a value $\beta^*$,
\begin{equation}\label{eq:kappa}
\gamma(\theta)=\E\big[\sigma_t^{-1}\partial_uA(u_t;c_t,\eta)\nabla_\beta\mu_t\big],\quad
J=\E\big[\sigma_t^{-2}\partial_u\psi(u_t;\eta)\nabla_\beta\mu_t\nabla_\beta\mu_t\tr\big],\quad
\kappa(\theta)=J^{-}\gamma(\theta),
\end{equation}
\begin{equation}\label{eq:xi}
\xi_t(\theta;\beta)=a_t(\theta;\beta)-\kappa(\theta)\tr s^\mu_t(\beta).
\end{equation}
For a constant location, $\kappa(\theta)=\E[\partial_\mu a_t]/\E[\partial_\mu s^\mu_t]$, derivatives taken with $\sigma_t$ fixed. The statistic is
\begin{equation}\label{eq:stat}
Z_T(\theta)=\frac{\sum_t\hat\xi_t(\theta)}{\{\sum_t\hat\xi_t(\theta)^2\}^{1/2}},\qquad \mathcal S_T=\sup_{\theta\in\Theta}Z_T(\theta),
\end{equation}
with $\hat\xi_t$ evaluated at $\hat\beta$ and at $\hat\kappa$ (sample averages of the derivatives in \eqref{eq:kappa}). Large values indicate downward impulses; $\sup_\theta\{-Z_T(\theta)\}$ tests for upward impulses.

\subsection{Assumptions}
\begin{assumption}[Null hypothesis]\label{A:null}
There is $\beta^*\in\operatorname{int}B$ such that, for every $t$, the law of $y_t-\mu_t(\beta^*)$ given $\F_{t-1}$ is symmetric about zero.
\end{assumption}
\noindent Nothing is assumed about the shape, scale, tails or volatility dynamics of the true conditional law; only the location of the fitted model must be right (see Lemma~\ref{lem:pseudo}).

\begin{assumption}[Estimator]\label{A:est}
$\sqrt T(\hat\beta-\beta^*)=O_p(1)$. Any $\sqrt T$-consistent estimator will do; it need not be the QMLE.
\end{assumption}

\begin{assumption}[Regularity]\label{A:reg}
(i) $\{y_t\}$ is strictly stationary and ergodic. (ii) On a neighborhood $N$ of $\beta^*$, $a_t(\theta;\beta)$ and $s^\mu_t(\beta)$ are twice continuously differentiable in $\beta$, with first and second derivatives dominated, uniformly in $\theta\in\Theta$, by $\F_t$-measurable variables with finite mean. (iii) $\theta\mapsto a_t(\theta;\beta^*)$ is Lipschitz with an $\F_t$-measurable constant $L_t$, $\E L_t^2<\infty$. (iv) $\Omega(\theta,\theta)=\E[\xi_t(\theta;\beta^*)^2]$ is bounded away from zero on $\Theta$, and $J$ is nonsingular on the span of $\nabla\mu_t$.
\end{assumption}

\subsection{Results}
\begin{lemma}[The test function is bounded for Student-$t$ $g$]\label{lem:bounded}
If $g$ is the Student-$t$ density with $\nu>0$ degrees of freedom, $\E Z^{\nu+1}<\infty$, and $c\le\bar c<\infty$, then $\sup_u|A(u;c,\nu)|\le2^{(\nu+3)/2}\,\E(1+\bar c^2Z^2/\nu)^{(\nu+1)/2}<\infty$. If $\sigma_t\ge\underline\sigma>0$ (as in GARCH with $\omega>0$), $a_t$ and $\xi_t$ are bounded and all their moments exist, whatever the tails of $y_t$.
\end{lemma}
\begin{proof}
For $w=u+v$, $1+u^2/\nu\le1+2w^2/\nu+2v^2/\nu\le2(1+v^2/\nu)(1+w^2/\nu)$, so $g(u+v)/g(u)\le\{2(1+v^2/\nu)\}^{(\nu+1)/2}$. Apply this with $v=\pm cZ$, take expectations and add. $\psi(u)=(\nu+1)u/(\nu+u^2)$ is bounded, so $s^\mu_t$ is bounded when $\sigma_t\ge\underline\sigma$.
\end{proof}

\begin{lemma}[Martingale differences and orthogonality]\label{lem:mds}
Under Assumption~\ref{A:null}, for every $\theta$, $\{\xi_t(\theta;\beta^*),\F_t\}$ is a martingale difference sequence and $\E[\nabla_\beta\xi_t(\theta;\beta^*)]=0$.
\end{lemma}
\begin{proof}
Given $\F_{t-1}$, $\xi_t(\theta;\beta^*)$ is an odd function of $u_t$ with $\F_{t-1}$-measurable coefficients, and $u_t$ is conditionally symmetric, so $\E[\xi_t\mid\F_{t-1}]=0$. For the derivative, $\partial u_t/\partial\beta=-\sigma_t^{-1}\nabla\mu_t-u_t\nabla\log\sigma_t$ and $\nabla c_t=-c_t\nabla\log\sigma_t$, so
\[
\nabla_\beta a_t=-\sigma_t^{-1}A_u\nabla\mu_t-u_tA_u\nabla\log\sigma_t-c_tA_c\nabla\log\sigma_t+A_\eta\nabla\eta .
\]
$A_u$ is even and $uA_u$, $A_c$, $A_\eta$ are odd in $u$; with $\F_{t-1}$-measurable multipliers the three odd terms have conditional mean zero, so $\E\nabla_\beta a_t=-\gamma(\theta)$. The same calculation gives $\E\nabla_\beta s^\mu_t=-J$. Hence $\E\nabla_\beta\xi_t=-\gamma+J\kappa=0$.
\end{proof}
\noindent The second statement is Neyman orthogonality: to first order, the statistic does not depend on how $\beta$ was estimated, because every non-location parameter enters $\xi_t$ through an even function of $u_t$ multiplied by an odd one.

\begin{theorem}[Null distribution]\label{thm:null}
Under Assumptions~\ref{A:null}--\ref{A:reg}, $T^{-1/2}\sum_{t=1}^T\hat\xi_t(\cdot)\Rightarrow\mathbb G(\cdot)$ in $\ell^\infty(\Theta)$, where $\mathbb G$ is a centered Gaussian process with covariance $\Omega(\theta,\theta')=\E[\xi_t(\theta;\beta^*)\xi_t(\theta';\beta^*)]$. Consequently $\mathcal S_T\Rightarrow\sup_\Theta\mathbb G(\theta)/\Omega(\theta,\theta)^{1/2}$. No long-run variance (HAC) estimator is needed, even when $y_t$ is strongly dependent through regimes or volatility.
\end{theorem}
\begin{proof}
A mean-value expansion around $\beta^*$ gives
\[
\begin{aligned}
T^{-1/2}\sum_t\xi_t(\theta;\hat\beta,\hat\kappa)={}&T^{-1/2}\sum_t\xi_t(\theta;\beta^*,\kappa)+\Big[T^{-1}\sum_t\nabla_\beta\xi_t(\theta;\bar\beta)\Big]\tr\sqrt T(\hat\beta-\beta^*)\\
&-(\hat\kappa-\kappa)\tr T^{-1/2}\sum_t s^\mu_t(\hat\beta).
\end{aligned}
\]
By ergodicity and Assumption~\ref{A:reg}(ii), the bracket converges to $\E\nabla_\beta\xi_t=0$ (Lemma~\ref{lem:mds}) uniformly in $\theta$; with Assumption~\ref{A:est} the second term is $o_p(1)$. The last term is $o_p(1)\cdot O_p(1)$ because $s^\mu_t(\beta^*)$ is itself an odd-in-$u_t$ martingale difference and $\hat\kappa\to\kappa$. The leading term is a sum of stationary ergodic martingale differences: finite-dimensional convergence follows from the martingale CLT \citep{hall1980}, and tightness from Assumption~\ref{A:reg}(iii) \citep{andrews1994,hansen1996}. $T^{-1}\sum\hat\xi_t^2$ is consistent for $\Omega$ by the same domination argument, and the martingale-difference property removes all cross-products.
\end{proof}

\paragraph{Critical values.} On a finite grid, $\sup_\Theta$ is the maximum of a Gaussian vector with the estimated correlation matrix of the $Z_T(\theta)$. For a continuum, the multiplier bootstrap $T^{-1/2}\sum_t\hat\xi_t(\theta)v_t$, $v_t$ i.i.d.\ $N(0,1)$ independent of the data, reproduces $\mathbb G$ conditionally on the data \citep{hansen1996}; it is valid because $\xi_t$ is a martingale difference. The supremum is needed because the jump scale $\theta$ is not identified under the null \citep{davies1977,davies1987,andrews1994ap,hansen1992}.

\paragraph{Local alternatives.} Let $X$ satisfy Assumptions~\ref{A:null}--\ref{A:reg} with a constant location, $\mu_t(\beta)=\mu$, which covers every model used in this paper (static Student-$t$, GARCH-, GJR- and APARCH-$t$ with constant mean, the centered HMM). The data are
\[
Y_t=X_t-B_tG_t,\qquad B_t\sim\text{Bernoulli}(\varepsilon_T),\quad\varepsilon_T=c/\sqrt T,\ c\ge0,
\]
with $\{B_t\}$, $\{G_t\}$ i.i.d., $G_t\ge0$, and $B$, $G$, $X$ mutually independent. Let $\F^Y_t=\sigma(Y_s,s\le t)$. Because $\sigma_t(\beta)$ is a function of the past, write $\xi_t(\theta;\beta;y;\mathcal P)$ for the test function evaluated at the current value $y$ with scale computed from a past $\mathcal P$. Define the conditional drift and covariance under the alternative,
\[
\begin{aligned}
\psi^{\varepsilon}_t(\theta)&=\E\big[\xi_t(\theta;\beta^*;X_t-G_t;Y_{1:t-1})\,\big|\,\F^Y_{t-1}\big],\\
v^{\varepsilon}_t(\theta,\theta')&=\E\big[\xi_t(\theta;\beta^*;X_t;Y_{1:t-1})\,\xi_t(\theta';\beta^*;X_t;Y_{1:t-1})\,\big|\,\F^Y_{t-1}\big],
\end{aligned}
\]
and their null counterparts $\psi^0_t$ and $v^0_t$, obtained by replacing $Y_{1:t-1}$ by $X_{1:t-1}$ and $\F^Y_{t-1}$ by $\F^X_{t-1}$.

\begin{assumption}[Local regularity]\label{A:local}
(i) $\xi_t$, $\nabla_\beta\xi_t$ and $s^\mu_t$ are bounded on $N\times\Theta$ (Lemma~\ref{lem:bounded} gives this for Student-$t$ $g$ when $\sigma_t\ge\underline\sigma>0$ and the derivatives of $\sigma_t$ are bounded). (ii) The past enters continuously: on the joint probability space of $(X,B,G)$, as $\varepsilon\to0$,
\[
\sup_t\E\sup_{\theta}\big|\psi^{\varepsilon}_t(\theta)-\psi^0_t(\theta)\big|\to0,\qquad
\sup_t\E\sup_{\theta,\theta'}\big|v^{\varepsilon}_t(\theta,\theta')-v^0_t(\theta,\theta')\big|\to0,
\]
and the same holds for the conditional means of $\nabla_\beta\xi_t$. (iii) $\sqrt T(\hat\beta-\beta^*)=O_p(1)$ under the local alternatives. (iv) The uniform law of large numbers for $T^{-1}\sum_t\nabla_\beta\xi_t(\theta;\beta)$ over $N\times\Theta$ holds under the local alternatives, and $\theta\mapsto\xi_t$ is Lipschitz with a bounded constant.
\end{assumption}
\noindent Condition (ii) says that a jump in the past changes the conditional behavior of later test terms by an amount that is small on average when jumps are rare. In GARCH-type models an impulse at $s$ moves $\sigma_t$ by an amount that decays geometrically in $t-s$, so (ii) follows from bounded continuous dependence of $\psi_t$ and $v_t$ on $\sigma_t$; in static models (ii) holds trivially because $\sigma_t$ does not depend on the past. Condition (iii) is implied by Theorem~\ref{thm:displacement}, under which the pseudo-true value moves by $O(\varepsilon_T)$, together with the usual $\sqrt T$-consistency of the estimator around it.

\begin{theorem}[Local power]\label{thm:local}
Under Assumptions~\ref{A:null}--\ref{A:reg} and \ref{A:local},
\[
\begin{aligned}
&T^{-1/2}\sum_{t=1}^T\hat\xi_t(\cdot)\Rightarrow\mathbb G(\cdot)+c\,\lambda(\cdot)\ \text{ in }\ell^\infty(\Theta),\\
&\lambda(\theta)=\E\,\psi^0_t(\theta)=\E_0\E_G\big[\xi_t(\theta;\beta^*;X_t-G;X_{1:t-1})\big],
\end{aligned}
\]
where $\mathbb G$ is the null limit of Theorem~\ref{thm:null}. Hence $Z_T(\theta)$ has asymptotic mean $c\,\lambda(\theta)/\Omega(\theta,\theta)^{1/2}$ and unit variance, and the local power of $\mathcal S_T$ is $\Pr\{\sup_\Theta[\mathbb G(\theta)+c\lambda(\theta)]/\Omega(\theta,\theta)^{1/2}>q_{1-\alpha}\}$. The drift is a single-impulse functional: an impulse at $t$ affects $\xi_t$ only, because $\sigma_t$ depends on the past alone and, for $u>t$, $Y_u=X_u$ remains conditionally symmetric about $\mu^*$.
\end{theorem}
\noindent The proof is in Appendix~\ref{app:proofs}. Under the local alternative, the conditional expectation of $\xi_t$ splits into a symmetric branch, with probability $1-\varepsilon_T$, and a jump branch, with probability $\varepsilon_T$; oddness removes the first, and the second gives the drift $c\lambda(\theta)$. A martingale central limit theorem for triangular arrays then gives the shifted Gaussian limit.

\begin{remark}[Finite-sample studentization]
Step~5 shows that, under the alternative, $T^{-1}\sum\hat\xi_t^2$ estimates $(1-\varepsilon)\Omega+\varepsilon\E[\xi_t^2|_{y_t=X_t-G}]$; the $O(\varepsilon)$ term vanishes asymptotically but matters in finite samples with large $\theta$ (Section~\ref{sec:mc}).
\end{remark}

\begin{remark}[Time-varying location]
If $\mu_t$ depends on the past (for example an AR mean), an impulse at $t$ also moves $\mu_u$ for $u>t$, and $Y_u-\mu_u(Y_{1:u-1})$ is no longer symmetric about zero given the past. The drift then acquires post-impulse terms, $\lambda(\theta)=\E_0\E_G\sum_{u\ge t}\{\xi_u(X^{[t,G]})-\xi_u(X)\}$, and Step~1 must be replaced by a mixing argument with square-summable impulse responses. We do not pursue this case.
\end{remark}

\begin{remark}[One functional for displacement and power]\label{rem:onefunctional}
$\lambda(\theta)$ has the same single-impulse form as the displacement in Theorem~\ref{thm:displacement}: it is that functional applied to the test function $\xi$ instead of the model's score. Theorem~\ref{thm:displacement} says how far an omitted jump moves each fitted model; Theorem~\ref{thm:local} says how detectable the same jump is. Because of orthogonality, the drift of $\hat\beta$ predicted by Theorem~\ref{thm:displacement} does not enter the local power at first order.
\end{remark}

\begin{remark}[The price of robustness]
The robust function is the downward-impulse score minus the upward-impulse score. Under a correct model this costs local power relative to the naive score; under a misspecified model the comparison can go either way. In the HMM-9 design (Table~\ref{tab:hmm9}) the robust test is more powerful than the naive one, and the naive test's size ranges from 0\% to 99.6\% across designs (Table~\ref{tab:full}).
\end{remark}

\begin{remark}[Two-sided jumps]\label{rem:twosided}
Jump models with separate upward and downward components \citep{kou2002,ramezani2007} fit into the same local analysis. Suppose that, in the setting of Theorem~\ref{thm:local}, downward impulses $-G_t$ occur with probability $\varepsilon^-_T=c^-/\sqrt T$ and upward impulses $+G_t$ with probability $\varepsilon^+_T=c^+/\sqrt T$, with the same size law. Because $X_t-\mu^*$ is conditionally symmetric and $\xi_t$ is odd about $\mu^*$, $\E_0\E_G\,\xi_t(X_t+G)=-\E_0\E_G\,\xi_t(X_t-G)$, so the drift becomes $(c^--c^+)\lambda(\theta)$. To first order the test measures the \emph{net} downward jump risk: equal upward and downward jumps cancel, as they should for a test robust to symmetric misspecification, and an excess of downward jumps is detected. When the two size laws differ, the drift is $c^-\lambda^-(\theta)-c^+\lambda^+(\theta)$, with $\lambda^\pm$ computed from the respective laws.
\end{remark}

\begin{proposition}[Consistency]\label{prop:cons}
Under a fixed alternative, if $\hat\beta\to_p\beta^\dagger$ and $\bar\lambda(\theta)=\E[\xi_t(\theta;\beta^\dagger)]>0$ for some $\theta\in\Theta$, then $\Pr(\mathcal S_T>q)\to1$ for every fixed $q$. For small $\varepsilon$, $\bar\lambda(\theta)=\varepsilon\lambda(\theta)+o(\varepsilon)$, so $\lambda(\theta)>0$ suffices.
% OPEN (A4): primitive condition for lambda(theta)>0.
\end{proposition}
\noindent The statistic is $\sqrt T\,\bar\lambda/\Omega^{1/2}+O_p(1)$. The condition can hold without jumps: any negatively skewed conditional law can give $\bar\lambda>0$. The test is therefore a test of one-sided, jump-like asymmetry against a symmetric null; Section~\ref{sec:b12} measures how often it rejects skewed innovations without jumps.

\begin{lemma}[When the pseudo-true location is the center]\label{lem:pseudo}
Suppose $\mu_t(\beta)=\mu$ is constant, $\sigma_t$ depends on the past only through $\{(y_s-\mu)^2\}_{s<t}$ or is constant, and the process is reflection invariant: $\{y_t-m\}$ and $\{-(y_t-m)\}$ have the same law. If the QMLE pseudo-true value is unique, its location is $\mu^*=m$, so Assumption~\ref{A:null} holds whenever $y_t$ is conditionally symmetric about $m$.
\end{lemma}
\begin{proof}
Reflection maps the expected log-likelihood at $(\mu,\beta_{-\mu})$ to that at $(2m-\mu,\beta_{-\mu})$; uniqueness forces $\mu^*=m$.
\end{proof}
\noindent Two-state Student-$t$ HMMs, GARCH and APARCH-type models without leverage, and symmetric contamination are reflection invariant. With leverage the assumption must be argued directly; in the application the rejections survive a GJR specification (Section~\ref{sec:daily}).
% OPEN (A5): leverage extension of the lemma.

\paragraph{Least-squares versus Jacobian projection.} A natural alternative projects $a_t$ on the location score with the least-squares coefficient $\kappa_{\rm OLS}=\sum a_ts^\mu_t/\sum(s^\mu_t)^2$. When $g$ is the true conditional law, the generalized information equality gives $\kappa_{\rm OLS}=\kappa$; under tail misspecification only the Jacobian ratio \eqref{eq:kappa} makes the statistic orthogonal, and since $\kappa_{\rm OLS}$ minimizes $\E(a_t-\kappa s^\mu_t)^2$ it can understate the variance. In a static model estimated by QMLE, $\sum_ts^\mu_t(\hat\beta)=0$, so the numerator is the same for both and only the studentization changes. In all our designs the two versions are numerically almost identical (decisions agree on 99.9\% of paths).

\section{After a rejection: when can the jumps be estimated?}\label{sec:after}
A rejection says that the data contain one-sided impulses relative to the fitted symmetric model. The risk manager must then choose between fitting a model with an explicit jump component and keeping the symmetric model while reporting its bias. Two jump models are natural. The first is a GARCH- or APARCH-$t$ model with Bernoulli--Gamma jumps ($-B_tG_t$). The second, which we call \emph{HMM-9}, is the two-state Student-$t$ hidden Markov model of Section~\ref{sec:setting} with Bernoulli--Gamma jumps in its emissions: in state $j$ the density of $Y_t$ is $(1-\varepsilon)f_j(y)+\varepsilon h_j(y;k,\theta)$, with $h_j(y;k,\theta)=\E f_j(y+G)$, so that the six parameters $\psi_0$ of the jump-free HMM (the HMM-6) are extended by the three jump parameters $(\varepsilon,k,\theta)$. In both models, the information about the shape and scale of the jumps comes only from the jumps themselves, and the choice depends on whether the sample contains enough of it.

Consider the per-observation law $p_\varepsilon=(1-\varepsilon)f+\varepsilon h_\varphi$, where $f$ is the jump-free density, $h_\varphi(y)=\E f(y+G)$ and $\varphi$ is a parameter of the jump-size law, either the scale $\theta$ or the shape $k$ of $G=\theta Z$, $Z\sim$ Gamma$(k,1)$. Its Fisher information about $\varphi$ is
\begin{equation}\label{eq:Iphi}
I_\varphi(\varepsilon)=\varepsilon^2\int\frac{(\partial_\varphi h_\varphi)^2}{(1-\varepsilon)f+\varepsilon h_\varphi}\,dy,
\qquad\text{and we write}\qquad C_q=\int\frac{(\partial_\varphi h_\varphi)^2}{q}\,dy .
\end{equation}

\begin{proposition}[Information about the jump law]\label{prop:jumpinfo}
For every $\varepsilon\in(0,1)$,
\[
\varepsilon^2\,C_{f+h}\;\le\;I_\varphi(\varepsilon)\;\le\;\min\Big\{\varepsilon\,C_h,\ \frac{\varepsilon^2}{1-\varepsilon}\,C_f\Big\}.
\]
If $f$ is a Student-$t$ density with $\nu>0$ degrees of freedom and scale $s>0$, then $C_f<\infty$ for $\varphi=\theta$ and $\varphi=k$. Consequently, with $T$ i.i.d.\ observations, the asymptotic variance of any regular estimator of $\varphi$, with or without the other parameters estimated, is at least
\[
\frac{1}{T\,I_\varphi(\varepsilon)}\;\ge\;\max\Big\{\frac{1}{T\varepsilon\,C_h},\ \frac{1-\varepsilon}{T\varepsilon^2\,C_f}\Big\}.
\]
\end{proposition}
\noindent The proof is in Appendix~\ref{app:proofs}.

\noindent The two upper bounds describe two regimes. When jumps are large relative to the scale of $f$, they land where $f$ is small, $\varepsilon C_h$ is the smaller bound, and the information is proportional to the number of jumps $T\varepsilon$. When jumps are moderate they overlap the tails of $f$; then $\varepsilon^2C_f$ binds and the information is only of order $\varepsilon^2$, the same order as the noncentrality of the test in Theorem~\ref{thm:local}. In that regime jumps can be \emph{detected} when $T\varepsilon^2$ is moderate, but their size and shape cannot be \emph{estimated} with useful precision. Table~\ref{tab:crb} evaluates the bounds for the designs of this paper. The information in \eqref{eq:Iphi} is exact for i.i.d.\ data; for a Markov-switching base it applies state by state and is a guide rather than a bound.

\begin{table}[ht]\centering\footnotesize\setlength{\tabcolsep}{4pt}
\caption{Fisher information about the jump scale $\theta$ and Cram\'er--Rao lower bounds for the standard deviations of $\hat\theta$ and $\hat k$. Base law: Student-$t$, $\nu=4$, unit scale; jumps $\theta\,$Gamma$(2,1)$. The smaller of the two upper bounds identifies the regime.}\label{tab:crb}
\begin{tabular}{lccclrcc}
\toprule
Design & $I_\theta$ & $\varepsilon C_h$ & $\varepsilon^2C_f$ & Regime & $T$ & sd$(\hat\theta)\ge$ & sd$(\hat k)\ge$\\
\midrule
HMM-9: $\varepsilon=0.03$, mean 6 & $8.9\times10^{-3}$ & $1.1\times10^{-2}$ & 0.19 & jump count & 1{,}000 & 0.33 & 0.38\\
 & & & & & 3{,}000 & 0.19 & 0.22\\
$\varepsilon=0.01$, mean 6 & $2.7\times10^{-3}$ & $3.7\times10^{-3}$ & $2.1\times10^{-2}$ & jump count & 4{,}800 & 0.28 & 0.32\\
$\varepsilon=0.02$, mean 1.8 & $3.9\times10^{-3}$ & $2.5\times10^{-2}$ & $4.8\times10^{-3}$ & overlap ($\varepsilon^2$) & 4{,}800 & 0.23 & 0.60\\
\bottomrule
\end{tabular}
\end{table}

\noindent The HMM-9 simulations agree. With $\varepsilon=0.03$ and mean jump 6, the bound for sd$(\hat k)$ is 0.38 at $n=1{,}000$ around a true value of 2, and Table~\ref{tab:hmm9info} shows $k$ at its search bound in 53 of 300 fits; at $n=3{,}000$ the bound falls to 0.22 and boundary hits fall to 6. For the moderate jumps of the main simulation design, the shape $k$ cannot be estimated better than $\pm0.60$ even with 4{,}800 observations, which is why a skewed innovation law and moderate jumps are hard to tell apart after a rejection (Section~\ref{sec:mc}). A practical rule follows: after a rejection, fit a jump model only when the estimated information $T\hat I_\theta$ makes the implied standard error small relative to $\hat\theta$; otherwise report the jump-induced bias of the symmetric model through Theorem~\ref{thm:displacement}. When $\varepsilon$ stays away from zero and the jump model is identified, standard interior HMM asymptotics apply \citep{douc2004}.

\begin{table}[ht]\centering\footnotesize\setlength{\tabcolsep}{4pt}
\caption{Estimating the jump parameters needs many jumps. HMM-9 fits to the negative-Gamma generator ($\varepsilon=0.03$, $k=2$, $\theta=3$); 300 paths per row for estimates and boundary hits, 500 paths for the Wald study ($n=1{,}000$ and $3{,}000$ only).}\label{tab:hmm9info}
\begin{tabular}{rrrrrrrr}
\toprule
$n$ & $n\varepsilon$ & Mean $\hat\varepsilon$ & Mean $\hat k$ & Mean $\hat\theta$ & $k$ at bound & Wald eligible & Joint 9-D coverage\\
\midrule
1{,}000 & 30 & 0.0321 & 3.207 & 2.670 & 53/300 & 57.4\% & 83.3\%\\
2{,}000 & 60 & 0.0321 & 2.498 & 2.984 & 18/300 & -- & --\\
3{,}000 & 90 & 0.0317 & 2.351 & 2.986 & 6/300 & 85.6\% & 89.7\%\\
\bottomrule
\end{tabular}
\end{table}

\subsection{A sequential procedure}\label{sec:procedure}
The results combine into a sequence of steps for a return series, each answering one question with its own tool.
\begin{enumerate}[itemsep=2pt]
\item \emph{Choose the base model.} The companion paper \citet{giornoA} shows that regime persistence can be estimated only when the effective sample $T\hat\kappa^2$, with $\hat\kappa$ the estimated contrast between regimes, is large (above about 100 in its designs). Use a two-state Student-$t$ HMM when it is; otherwise a GARCH- or APARCH-$t$ model.
\item \emph{Test for downside impulses} with the robust test of Section~\ref{sec:test}, using the base model as the null. If the test does not reject, keep the symmetric base model.
\item \emph{Pre-check whether the jumps can be described.} The location of the maximum of $Z_T(\theta)$ gives a rough jump scale $\tilde\theta$, and Theorem~\ref{thm:local} gives a rough intensity $\tilde\varepsilon\approx Z_T(\tilde\theta)\,\hat\Omega(\tilde\theta,\tilde\theta)^{1/2}/\{\sqrt T\hat\lambda(\tilde\theta)\}$, both computed from the fitted base model. Plugging them into \eqref{eq:Iphi} gives $n_J=T\,I_\theta\,\theta^2$, the inverse squared relative Cram\'er--Rao error of $\hat\theta$. For a model with volatility dynamics, the jump enters in standardized units $\theta/\sigma_t$, so $n_J=\theta^2\sum_tI_\theta(\varepsilon,\theta/\sigma_t)/\sigma_t^2$: jumps are easier to describe in calm periods. The pre-check is a screen, not an estimate, because $\tilde\varepsilon$ also absorbs skewness.
\item \emph{If $n_J$ is small} ($n_J$ well below 16, a relative error well above 25\%), the data are in the overlap regime. Do not fit a jump model. Keep the base model and report the jump-induced displacement of its parameters and of its VaR (Theorem~\ref{thm:displacement}, Corollaries~\ref{cor:var}--\ref{cor:regret}); a skewed innovation law, such as the Fern\'andez--Steel skew-$t$, is a practical alternative for forecasting. \emph{If $n_J$ is large}, the data are in the jump-count regime: fit the extended model, the HMM-9 when step~1 chose regimes and a GARCH- or APARCH-$t$ model with Bernoulli--Gamma jumps otherwise. Table~\ref{tab:hmm9info} shows that the HMM-9 needs many jumps (about 90 in our design) for reliable inference, so in most samples the second option is the practical one.
\item \emph{Confirm $n_J$ at the fitted estimates} before reporting jump parameters and their standard errors; near the boundary, prefer a bootstrap. The existence of jumps is established by the test of step~2, never by a $t$-statistic for $\varepsilon$, whose null distribution is nonstandard (Section~\ref{sec:test}).
\end{enumerate}

\section{Monte Carlo evidence}\label{sec:mc}
\subsection{Designs}
Two independent sets of designs are used. The \emph{HMM-6 manifest} has eight two-state Student-$t$ HMM designs with 500 paths each ($T=6{,}000$; the first 4{,}800 observations are the estimation sample), regenerated exactly from \texttt{DGP\_MANIFEST.csv} (all 4{,}000 seeds reproduced from \texttt{SeedSequence(20261028)}; returns identical to the archived generator). The jump arm adds downward impulses with $\varepsilon=0.02$ and size $0.9\,$Gamma$(2,1)$ (mean 1.8) from a separate seeded stream. The \emph{HMM-9 design} (archived from an earlier study) has nine cells of 300 paths: a clean two-state Student-$t$ HMM ($p_{00}=0.985$, $p_{11}=0.960$, $q=(1.12,1.42)$, $s=(0.80,1.25)$), the same HMM with symmetric $N(0,5^2)$ contamination at rate 0.03, and the same HMM with negative Gamma$(2,3)$ jumps at rate 0.03, for $n=1{,}000$, $2{,}000$, $3{,}000$.

\subsection{Where the jumps go: cross-model displacement}\label{sec:cross}
In the full HMM-6 design with $k=2$, the predicted derivatives $\bar H_M^{-1}\bar b_M$ use 20{,}000 single impulses on a 400{,}000-observation path; the simulated derivatives are finite differences of pseudo-true fits on the same path with common random numbers. At $\varepsilon=0.01$, 28 of 32 components lie within two combined standard errors (Table~\ref{tab:cross}). The location shifts are 2--5\% larger than predicted, and the APARCH tail index at $\theta=0.9$ curves in $\varepsilon$; both are effects of higher order that the first-order theory does not capture. The single-impulse calculation therefore predicts the direction and size of the displacement in every model well, without fitting any model to jump paths.
% OPEN (A6, B5). In the 240-path jump pilot, median paired parameter shifts agree with the large-sample shifts at $\varepsilon=0.02$, for example Student-$t$ location $-0.017$ against $-0.016$, APARCH $\gamma$ $-0.081$ against $-0.080$, and HMM $p_{11}$ $-0.022$ against $-0.018$.

\begin{table}[ht]\centering\footnotesize\setlength{\tabcolsep}{4pt}
\caption{Predicted and finite-difference (FD) derivatives of the pseudo-true parameters with respect to the jump intensity $\varepsilon$. Data-generating process: full HMM-6 design, jumps $\theta\,$Gamma$(2,1)$.}\label{tab:cross}
\begin{tabular}{lllrrrrr}
\toprule
$\theta$ & Model & Param. & Predicted (s.e.) & FD $\varepsilon{=}.0025$ & FD $.005$ & FD $.01$ & $z$ at $.01$\\
\midrule
0.9 & Student-$t$ & $\mu$ & $-0.828$ (0.006) & $-0.848$ & $-0.819$ & $-0.857$ & $-2.08$\\
0.9 & Student-$t$ & $s$ & 0.322 (0.019) & 0.393 & 0.346 & 0.378 & 1.21\\
0.9 & Student-$t$ & $\nu$ & $-17.218$ (0.391) & $-16.582$ & $-17.073$ & $-15.837$ & 1.44\\
0.9 & APARCH-$t$ & $\mu$ & $-0.797$ (0.007) & $-0.815$ & $-0.796$ & $-0.830$ & $-1.93$\\
0.9 & APARCH-$t$ & $\omega$ & 0.518 (0.216) & 0.618 & 0.457 & 0.379 & $-0.26$\\
0.9 & APARCH-$t$ & $\alpha$ & $-0.288$ (0.069) & $-0.256$ & $-0.262$ & $-0.376$ & $-0.52$\\
0.9 & APARCH-$t$ & $\gamma$ & $-3.905$ (0.098) & $-4.521$ & $-4.476$ & $-3.958$ & $-0.22$\\
0.9 & APARCH-$t$ & $\beta$ & 1.703 (0.287) & 1.449 & 1.601 & 1.578 & $-0.18$\\
0.9 & APARCH-$t$ & $\delta$ & $-20.279$ (0.949) & $-18.240$ & $-18.725$ & $-16.161$ & 1.77\\
0.9 & APARCH-$t$ & $\nu$ & $-35.806$ (0.588) & $-33.147$ & $-33.451$ & $-30.819$ & 3.46\\
0.9 & HMM-6 & $p_{00}$ & $-0.450$ (0.014) & $-0.435$ & $-0.465$ & $-0.439$ & 0.34\\
0.9 & HMM-6 & $p_{11}$ & $-1.056$ (0.043) & $-1.071$ & $-1.124$ & $-0.998$ & 0.56\\
0.9 & HMM-6 & $\nu_0$ & $-438.8$ (22.4) & $-453.0$ & $-416.0$ & $-382.3$ & 1.03\\
0.9 & HMM-6 & $\nu_1$ & $-29.890$ (2.545) & $-27.012$ & $-25.461$ & $-24.211$ & 0.91\\
0.9 & HMM-6 & $s_0$ & $-0.397$ (0.034) & $-0.356$ & $-0.410$ & $-0.394$ & 0.03\\
0.9 & HMM-6 & $s_1$ & $-0.862$ (0.257) & $-0.467$ & $-0.412$ & $-0.511$ & 0.56\\
\addlinespace
0.4 & Student-$t$ & $\mu$ & $-0.611$ (0.004) & $-0.637$ & $-0.617$ & $-0.642$ & $-2.99$\\
0.4 & Student-$t$ & $s$ & 0.352 (0.011) & 0.378 & 0.374 & 0.376 & 0.93\\
0.4 & Student-$t$ & $\nu$ & $-0.385$ (0.173) & $-0.193$ & $-0.394$ & $-0.015$ & 0.87\\
0.4 & APARCH-$t$ & $\mu$ & $-0.612$ (0.005) & $-0.631$ & $-0.617$ & $-0.641$ & $-2.63$\\
0.4 & APARCH-$t$ & $\omega$ & 0.147 (0.153) & 0.141 & 0.135 & 0.093 & $-0.14$\\
0.4 & APARCH-$t$ & $\alpha$ & $-0.149$ (0.047) & $-0.106$ & $-0.086$ & $-0.150$ & $-0.01$\\
0.4 & APARCH-$t$ & $\gamma$ & $-0.505$ (0.046) & $-0.802$ & $-0.838$ & $-0.634$ & $-1.15$\\
0.4 & APARCH-$t$ & $\beta$ & 0.263 (0.202) & 0.209 & 0.235 & 0.289 & 0.05\\
0.4 & APARCH-$t$ & $\delta$ & $-1.440$ (0.520) & $-1.010$ & $-1.667$ & $-1.125$ & 0.25\\
0.4 & APARCH-$t$ & $\nu$ & $-2.708$ (0.286) & $-1.761$ & $-2.177$ & $-1.733$ & 1.39\\
0.4 & HMM-6 & $p_{00}$ & $-0.015$ (0.007) & $-0.033$ & $-0.022$ & $-0.007$ & 0.51\\
0.4 & HMM-6 & $p_{11}$ & $-0.052$ (0.019) & $-0.113$ & $-0.088$ & $-0.010$ & 0.89\\
0.4 & HMM-6 & $\nu_0$ & $-167.2$ (13.3) & $-132.0$ & $-165.3$ & $-141.1$ & 0.80\\
0.4 & HMM-6 & $\nu_1$ & 1.286 (0.827) & 3.094 & 3.398 & 1.477 & 0.09\\
0.4 & HMM-6 & $s_0$ & 0.132 (0.022) & 0.179 & 0.133 & 0.175 & 0.81\\
0.4 & HMM-6 & $s_1$ & 0.333 (0.102) & 0.481 & 0.528 & 0.297 & $-0.14$\\
\bottomrule
\end{tabular}
\end{table}

\subsection{What the jumps cost: value-at-risk}\label{sec:var}
All 4{,}000 manifest paths are used, with 4{,}800 estimation and 1{,}200 evaluation observations. Each path has three versions on the same regime path: no jumps, and Bernoulli--Gamma jumps ($-B_tG_t$, size $0.9\,$Gamma$(2,1)$, mean 1.8) at rates $\varepsilon=0.01$ and $0.02$, the first a subset of the second's jump dates. Four models are fitted without a jump component and filtered through the evaluation sample with fixed parameters: a static Student-$t$, GARCH(1,1)-$t$, APARCH(1,1)-$t$ \citep{ding1993} and the HMM-6. The benchmark is the true conditional value-at-risk, computed from the data-generating HMM including its jump component by the exact forward filter. Every HMM-6 fit converged.

\begin{table}[ht]\centering\footnotesize\setlength{\tabcolsep}{4pt}
\caption{Value-at-risk at 99\% and expected shortfall at 97.5\%, 4{,}000 paths per cell. Shortfall is the mean of (true $-$ model)/true; positive means under-capitalised. Basel zone from exceptions in the last 250 evaluation days \citep{basel1996}; capital is the Basel charge relative to the charge implied by the true VaR. The static Student-$t$ ignores volatility dynamics and is conservative on average, so its shortfalls are not comparable with the conditional models.}\label{tab:var}
\begin{tabular}{llrrrrrr}
\toprule
Model & $\varepsilon$ & Breach & Kupiec rej. & VaR shortf. & ES shortf. & Yellow & Capital\\
 & & \% & \% & \% & \% & \% & ratio\\
\midrule
GARCH-$t$ & 0 & 1.06 & 8.6 & $-0.2$ & $-0.1$ & 13.5 & 1.003\\
 & 0.01 & 1.17 & 11.7 & 3.0 & 3.4 & 17.8 & 0.980\\
 & 0.02 & 1.28 & 17.2 & 6.0 & 6.1 & 21.8 & 0.959\\
\addlinespace
APARCH-$t$ & 0 & 1.04 & 7.1 & $-0.9$ & $-1.0$ & 12.5 & 1.008\\
 & 0.01 & 1.15 & 9.8 & 2.4 & 2.5 & 17.1 & 0.984\\
 & 0.02 & 1.25 & 14.5 & 5.4 & 5.1 & 20.6 & 0.963\\
\addlinespace
HMM-6 & 0 & 1.01 & 5.6 & $-0.05$ & $-0.02$ & 11.1 & 1.001\\
 & 0.01 & 1.15 & 8.6 & 3.6 & 3.8 & 17.0 & 0.976\\
 & 0.02 & 1.27 & 15.6 & 6.7 & 6.7 & 21.1 & 0.953\\
\addlinespace
Student-$t$ & 0 & 1.03 & 8.0 & $-8.6$ & $-11.1$ & 12.8 & 1.051\\
 & 0.01 & 1.12 & 9.4 & $-3.9$ & $-6.2$ & 16.1 & 1.021\\
 & 0.02 & 1.21 & 13.4 & 0.0 & $-2.5$ & 19.4 & 0.994\\
\bottomrule
\end{tabular}
\end{table}

Without jumps the conditional models are almost exactly calibrated: the HMM-6, which is correctly specified there, has a breach rate of 1.01\% and a VaR shortfall of $-0.05\%$. With jumps at 2\% intensity, every conditional model leaves 99\% VaR 5.4--6.7\% too low and 97.5\% expected shortfall 5.1--6.7\% too low; breach rates rise to 1.25--1.28\%, Kupiec rejections double or more (from 5.6--8.6\% to 14.5--17.2\%), Basel yellow-zone outcomes rise from 11--14\% to 21--22\%, and the regulatory capital falls 4--5\% short of what the true VaR requires. The results agree with the 240-path pilot (VaR shortfall 6.0--6.9\%, yellow zone 20--21\%). With 4{,}800 estimation and 1{,}200 evaluation observations, the ratio of the two samples is 0.25, so the Kupiec rates carry the estimation-risk distortion described by \citet{escanciano2010}; the correctly specified HMM-6 rejects 5.6\% of jump-free paths, close to the nominal 5\%, which suggests that the distortion is small here. The shortfalls are largest in the designs with large variance contrasts (GARCH-$t$: 11.9\% in the calibration-anchor design), where the models absorb jumps into their scale parameters.

\paragraph{First-order error, second-order regret.} Corollaries~\ref{cor:var} and~\ref{cor:regret} predict that the VaR error is first order in $\varepsilon$ and the quantile-loss regret of a correctly specified model second order. Doubling $\varepsilon$ from 0.01 to 0.02 should therefore double the increase in VaR shortfall and multiply the increase in regret by four. The shortfall increases by factors of 1.85--1.93 for all four models. For the HMM-6, the increase in expected pinball regret (computed under the true conditional law) grows by a factor of 3.37 at 95\%, 3.49 at 99\% and 3.31 at 99.5\%, with path-bootstrap 95\% intervals of $[3.30,3.43]$, $[3.43,3.56]$ and $[3.25,3.36]$. The ratio is close to the second-order value of 4 and far from the first-order value of 2. It stays below 4 because the estimated HMM-6 is not exactly correct at $\varepsilon=0$: its sampling error adds a first-order cross term $\varepsilon\,\E[f_t\Delta_t(0)\dot\Delta_t]$ that vanishes as $T$ grows. For models that are already misspecified at $\varepsilon=0$ the incremental regret is first order and can have either sign: the static Student-$t$, whose VaR is too conservative without jumps, has \emph{lower} regret when jumps move the true quantile towards its forecast.

\subsection{Detecting the jumps: size and power on the HMM-6 manifest}\label{sec:full}
The test uses the training samples, the grid $\theta\in\{0.3,0.6,0.9,1.5,2.5\}$ and two symmetric null models, a static Student-$t$ and a GARCH(1,1)-$t$, both misspecified for HMM data.

\begin{table}[ht]\centering\footnotesize\setlength{\tabcolsep}{4pt}
\caption{Rejection rates (\%) at the 5\% level, 500 paths per design. ``Robust'': orthogonalized statistic of Theorem~\ref{thm:null}; ``Naive'': downward-only score. Monte Carlo s.e.\ about 1.0 point per cell at 5\% and 0.3 points for the pooled size.}\label{tab:full}
\begin{tabular}{l rr rr rr rr}
\toprule
 & \multicolumn{4}{c}{Without jumps (size)} & \multicolumn{4}{c}{With jumps (power)}\\
\cmidrule(lr){2-5}\cmidrule(lr){6-9}
 & \multicolumn{2}{c}{Robust} & \multicolumn{2}{c}{Naive} & \multicolumn{2}{c}{Robust} & \multicolumn{2}{c}{Naive}\\
Design & $t$ & GARCH & $t$ & GARCH & $t$ & GARCH & $t$ & GARCH\\
\midrule
Calibration anchor & 5.6 & 4.0 & 99.6 & 81.0 & 40.0 & 62.8 & 100.0 & 99.6\\
Full HMM-6 & 4.4 & 4.4 & 0.0 & 0.0 & 64.2 & 77.0 & 31.2 & 24.8\\
Scale only & 5.2 & 4.8 & 0.8 & 0.0 & 65.4 & 72.8 & 49.8 & 37.6\\
Tail only & 4.0 & 4.2 & 0.0 & 0.0 & 90.6 & 90.2 & 32.6 & 33.8\\
Weak tail & 2.6 & 2.6 & 0.0 & 0.0 & 89.6 & 88.6 & 29.6 & 28.0\\
Weak persistence & 4.8 & 4.8 & 0.0 & 0.0 & 80.8 & 82.0 & 23.4 & 25.6\\
IID mixture & 3.8 & 3.8 & 0.0 & 0.0 & 87.8 & 87.6 & 32.6 & 32.0\\
Single Student-$t$ & 5.6 & 5.6 & 0.0 & 0.4 & 84.8 & 84.6 & 29.4 & 28.8\\
\midrule
All designs & 4.5 & 4.3 & 12.6 & 10.2 & 75.4 & 80.7 & 41.1 & 38.8\\
\bottomrule
\end{tabular}
\end{table}

The robust test keeps its size in every design (2.6--5.6\%; pooled 4.5\% and 4.3\%). The least-squares projection gives the same figures to within 0.2 points. The naive score is unusable: it rejects 99.6\% (Student-$t$) and 81.0\% (GARCH-$t$) of jump-free paths in the calibration-anchor design and almost never elsewhere, so its size depends entirely on how the tails are misspecified; its apparent power of 100\% in that design is meaningless.

The robust test detects jumps in 75.4\% (Student-$t$) and 80.7\% (GARCH-$t$) of paths, and its power is driven by whether the regimes differ in variance. In the five designs where both regimes have unit variance, power is 80.8--90.6\%. In the three where the variances differ (scale only: 0.56 against 3.63; full HMM-6: 0.55 against 3.42; calibration anchor: 0.30 against 5.83), it falls to 64--65\% and 40\% with the Student-$t$ null. This is the trade-off of Remark~\ref{rem:onefunctional}: the larger the variance contrast, the more an omitted jump can be absorbed into the scale and tail parameters (Theorem~\ref{thm:displacement}), and the harder it is to see. The GARCH-$t$ null, which partly tracks the scale changes, recovers part of the loss (72.8\%, 77.0\% and 62.8\%).

\subsection{Asymmetry versus jumps, and competing tests}\label{sec:b12}
A rejection means one-sided asymmetry, which jumps produce but so does a skewed innovation law (Proposition~\ref{prop:cons}). To measure both, each manifest path is given five versions built from the same regime path: the original (\emph{sym}); moderate jumps as above (\emph{jump}: $\varepsilon=0.02$, mean 1.8); large jumps (\emph{jumpbig}: $\varepsilon=0.01$, size $3\,$Gamma$(2,1)$, mean 6, on a subset of the same jump dates); and two skewed versions without jumps (\emph{skew90}, \emph{skew95}), in which each regime's Student-$t$ emission is replaced by its Fern\'andez--Steel skewing \citep{fernandez1998} with $\gamma=0.90$ or $0.95$, keeping $|X_t|$ and the regime path, and centered so that the conditional mean is unchanged. Results cover 3{,}924 of the 4{,}000 paths (490--491 per design).
% OPEN: complete the remaining 76 paths.

Two competitors are computed on the GARCH(1,1)-$t$ standardized residuals: a skewness test in the spirit of \citet{bainh2005}, based on $z_t^3-3z_t$ with a Newey--West variance and one-sided against negative skewness, and an omnibus symmetry statistic in the spirit of \citet{bainh2001}, $\sqrt T\sup_x|F_T(-x)-\{1-F_T(x^-)\}|$, with sign-flip critical values in which every replicate re-estimates the location. The re-estimation matters: the fitted location makes the residuals look more symmetric than they are, and without it the sign-flip test almost never rejects, even under the alternative (in our designs, null $p$-values had median 0.86). This estimation effect is what the martingale transformation of \citet{bainh2001} removes. The robust test needs no such correction, because its Jacobian projection makes it orthogonal to the estimated location (Lemma~\ref{lem:mds}).

\begin{table}[ht]\centering\footnotesize\setlength{\tabcolsep}{4pt}
\caption{Rejection rates (\%) at the 5\% level, pooled over the eight designs (3{,}924 paths; 4{,}000 for the symmetry test). The naive score's rates are not comparable because its size is 10.1\% pooled and 80.2\% in the calibration-anchor design; the symmetry test's size is 20.4\% in the calibration-anchor design and 8.8\% in the full HMM-6 design.}\label{tab:b12}
\begin{tabular}{lrrrrr}
\toprule
Data-generating process & Robust, $t$ & Robust, GARCH-$t$ & Skewness test & Symmetry test & Naive score\\
\midrule
No jumps (size) & 4.5 & 4.3 & 4.9 & 7.5 & 10.1\\
Moderate jumps (mean 1.8) & 75.4 & 80.8 & 64.9 & 16.8 & 38.8\\
Large jumps (mean 6) & 96.1 & 98.6 & 96.0 & 25.0 & 99.4\\
Skewed innovations ($\gamma=0.95$) & 55.6 & 56.8 & 37.3 & 36.0 & 15.6\\
Skewed innovations ($\gamma=0.90$) & 92.3 & 93.5 & 71.4 & 91.4 & 63.6\\
\bottomrule
\end{tabular}
\end{table}

Pooled over the designs, the robust test is the most powerful valid test in every version: against moderate jumps it rejects 80.8\% of paths against 64.9\% for the skewness test, and against the milder skewing 56.8\% against 37.3\%. The skewness test keeps its nominal size here, but it relies on sixth moments of the standardized residuals and loses most power where the regimes differ in variance (moderate jumps in the calibration-anchor design: 15.7\% against 63.1\% for the robust test). As Proposition~\ref{prop:cons} anticipated, the robust test also rejects skewed innovations without jumps (93.5\% at $\gamma=0.90$).

The omnibus symmetry test is the weakest against jumps. On the six designs where all three valid tests are correctly sized (size 4.3--5.1\%), the robust test, the skewness test and the symmetry test reject 84.4\%, 77.2\% and 13.4\% of paths with moderate jumps; 99.9\%, 98.4\% and 23.5\% with large jumps; 64.9\%, 44.9\% and 35.5\% with mild skewing; and 99.5\%, 83.4\% and 93.4\% with stronger skewing. The symmetry statistic compares the two halves of the whole distribution, so a handful of large downward observations barely moves it; it is competitive only against an asymmetry that spreads through the body of the distribution. In the two designs where it is oversized, its size-adjusted power against jumps is 12--20\%.

\paragraph{Telling skewness from jumps after a rejection.} On every path where the robust GARCH-$t$ test rejects, the GARCH volatility filter is held fixed and two innovation laws are fitted to the standardized residuals: a Fern\'andez--Steel skew-$t$ and a Student-$t$ with Bernoulli--Gamma downward jumps ($-B_tG_t$, shape $k=2$). They are compared by the \citet{vuong1989} statistic with a Newey--West variance and a Schwarz correction, choosing a model when $|V|>1.96$.

\begin{table}[ht]\centering\footnotesize\setlength{\tabcolsep}{4pt}
\caption{Model chosen by the Vuong comparison among paths where the robust GARCH-$t$ test rejects (\% of rejected paths).}\label{tab:vuong}
\begin{tabular}{lrrrr}
\toprule
Data-generating process & Rejected paths & Jumps & Skew-$t$ & Undecided\\
\midrule
Moderate jumps (mean 1.8) & 3{,}172 & 0.3 & 10.4 & 89.3\\
Large jumps (mean 6) & 3{,}870 & 76.4 & 0.8 & 22.8\\
Skewed innovations ($\gamma=0.95$) & 2{,}229 & 0.0 & 72.5 & 27.5\\
Skewed innovations ($\gamma=0.90$) & 3{,}669 & 0.4 & 74.2 & 25.4\\
\bottomrule
\end{tabular}
\end{table}

The pattern is the one predicted by Proposition~\ref{prop:jumpinfo}. Large jumps are identified as jumps in 76.4\% of rejected paths, and in 98--100\% in the five designs whose regimes have equal variance; the exceptions are the designs with large variance contrasts (full HMM-6 73.9\%, scale only 31.6\%, calibration anchor 0.5\%), where the volatility filter absorbs part of the jumps. Skewed innovations are almost never mistaken for jumps (at most 0.4\%) and are recognized as skewness in 72--74\% of rejected paths. Moderate jumps, by contrast, cannot be classified: 89.3\% of comparisons are undecided and 10.4\% wrongly favor skewness. These are exactly the jumps for which Table~\ref{tab:crb} shows that the shape of the jump law cannot be estimated with useful precision. Moderate jumps can be detected, and their effect on fitted models and on value-at-risk measured through Theorem~\ref{thm:displacement}, but they cannot be told apart from a skewed innovation law on samples of this size.

\section{Application to U.S.\ and Brazilian equity returns}\label{sec:data}
\subsection{The procedure on long daily samples}\label{sec:daily}
We apply the procedure of Section~\ref{sec:procedure} to three long daily series of adjusted closing prices from Yahoo Finance: the SPY exchange-traded fund from its first trading day (29 January 1993), the S\&P 500 index from 1990, and the Ibovespa from 1995, after the Real Plan, so that the hyperinflation years and index redenominations, which would look like enormous jumps, are excluded. Returns are $100\log(P_t/P_{t-1})$, and the samples end on 3 November 2026. The largest moves are genuine crisis days (October 2008 and March 2020 in the U.S.\ series; the January 1999 devaluation and the 1997--98 crises for the Ibovespa). For step~5 the models are estimated on the first 80\% of each sample and evaluated on the last 20\% (roughly 2020--2026) with fixed parameters.

\begin{table}[ht]\centering\footnotesize\setlength{\tabcolsep}{4pt}
\caption{The practical procedure on long daily samples. Step~1 uses GARCH(1,1)-$t$ standardized returns. Step~2 $p$-values are for downward impulses; the test for upward impulses rejects in none of the samples. Step~5: out-of-sample breach rates of 99\% VaR and the gain in 95\% pinball loss over GARCH-$t$, with Newey--West $t$-statistics.}\label{tab:daily}
\begin{tabular}{llccc}
\toprule
Step & & SPY & S\&P 500 & Ibovespa\\
 & & 1993--2026 & 1990--2026 & 1995--2026\\
\midrule
 & Returns & 8{,}498 & 9{,}277 & 7{,}885\\
1 & $T\hat\kappa^2$; persistence $Z$ & 32; $-3.7$ & 48; $-5.4$ & 1.7; $-4.3$\\
 & Base model & GARCH-$t$ & GARCH-$t$ & GARCH-$t$\\
2 & Robust test $p$, GARCH / GJR & $<0.001$ / $<0.001$ & $<0.001$ / $<0.001$ & 0.004 / 0.040\\
3 & Pre-check $\tilde\varepsilon$; $n_J$ & 0.12; flagged & 0.11; flagged & 0.034; 56\\
4 & Jump model: $\hat\varepsilon$ & 0.10 (bound) & 0.10 (bound) & 0.10 (bound)\\
 & Skew-$t$ $\hat\gamma$; Vuong & 0.885; $-1.0$ & 0.897; $-1.2$ & 0.929; $+2.1$\\
5 & 99\% breaches \%: GARCH / jumps / skew & 1.5 / 1.2 / 1.4 & 1.8 / 1.3 / 1.3 & 0.7 / 0.5 / 0.5\\
 & 95\% pinball gain, jumps & $+1.8\%$ (3.5) & $+1.5\%$ (3.9) & $-1.4\%$ ($-3.0$)\\
 & 95\% pinball gain, skew-$t$ & $+1.9\%$ (1.5) & $+1.5\%$ (1.5) & $-0.8\%$ ($-2.1$)\\
\bottomrule
\end{tabular}
\end{table}

\emph{Step 1.} On GARCH-standardized returns, the effective sample for regime persistence is 32--48 for the U.S.\ series and 1.7 for the Ibovespa, and the persistence test of \citet{giornoA} is negative in all three. Volatility clustering, which the GARCH filter removes, is the only persistent component, so the base model is GARCH-$t$ throughout. (On raw returns the same test is strongly positive; that is volatility clustering, not regimes, which is why step~1 standardises first.)

\emph{Step 2.} The robust test rejects conditional symmetry for downward impulses in all three series, and never for upward impulses. The rejection survives a leverage specification. For the U.S.\ series the statistic peaks at $Z_T\approx7$.

\emph{Steps 3--4.} The asymmetry does not behave like rare, large jumps. For both U.S.\ series the pre-check implies a jump intensity above 10\% per day, which is not interpretable as rare impulses. In all three series the joint maximum-likelihood fit of GARCH-$t$ with Bernoulli--Gamma jumps ($-B_tG_t$) pushes the intensity to its 10\% bound, with mean jump sizes of 0.7--1.0\%, below one daily standard deviation, and a singular Hessian: the likelihood turns the asymmetry into frequent moderate impulses, which amounts to a skewed innovation law. The Fern\'andez--Steel skew-$t$ fits with $\hat\gamma=0.885$--$0.929$. The Vuong comparison cannot separate the two for the U.S.\ series and favors the skew-$t$ for the Ibovespa, where the pre-check had passed ($n_J=56$). These are the overlap conditions of Proposition~\ref{prop:jumpinfo} and Section~\ref{sec:b12}: the asymmetry can be detected, but it cannot be attributed to jumps rather than skewness on these samples, and jump parameters should not be reported.

\emph{Step 5.} Accounting for the asymmetry matters for risk in the U.S.\ series. Out of sample, GARCH-$t$ under-covers: at 99\% its breach rate is 1.5\% for SPY and 1.8\% for the S\&P 500 (Kupiec $p=0.042$ and $0.002$), and at 95\% it is 6.6--6.9\% ($p\le0.002$). Either asymmetric model restores calibration (99\% breaches 1.2--1.4\%) and lowers the 95\% pinball loss by 1.5--1.9\%, significantly for the jump specification ($t=3.5$ and $3.9$); at 99\% the gains are of similar size but not significant. For the Ibovespa, whose GARCH-$t$ VaR was already conservative over 2020--2026, both asymmetric models are worse out of sample ($t$ between $-2.1$ and $-3.0$). A single 80/20 split is a limited evaluation design, and the Kupiec $p$-values are not corrected for estimation risk \citep{escanciano2010}, so they should be read as indicative; rolling re-estimation is a natural extension.

\subsection{Minute data and sampling frequency}\label{sec:minute}
The test uses the grid $\theta\in\{0.3,0.6,0.9,1.5,2.5\}\times\hat s$, with $\hat s$ the fitted Student-$t$ scale, and three symmetric nulls: static Student-$t$, GARCH(1,1)-$t$ and GJR-GARCH(1,1)-$t$. Minute samples use intraday returns only; overnight returns, which have a different scale, are excluded. Table~\ref{tab:data} reports $p$-values for downward impulses; the test for upward impulses rejects in none of these samples.

\begin{table}[ht]\centering\footnotesize\setlength{\tabcolsep}{4pt}
\caption{$p$-values of the robust test for downward impulses. SPY daily: adjusted closes, 2020--2026. Minute data: SPY January--October 2026 and WIN (Ibovespa mini-futures) February--October 2026; first 80\%/last 20\% as in the VaR study.}\label{tab:data}
\begin{tabular}{lrrrr}
\toprule
Series & $T$ & Student-$t$ & GARCH-$t$ & GJR-$t$\\
\midrule
SPY daily 2020--2026 & 1{,}714 & 0.005 & $<0.001$ & $<0.001$\\
SPY daily 2021--2026 (without 2020) & 1{,}462 & 0.010 & $<0.001$ & $<0.001$\\
\addlinespace
SPY 1-min, full & 76{,}080 & 0.016 & 0.060 & 0.081\\
SPY 1-min, first 80\% & 60{,}869 & 0.045 & 0.600 & 0.192\\
SPY 1-min, last 20\% & 15{,}211 & 0.100 & 0.016 & 0.023\\
SPY 5-min, full & 15{,}588 & 0.076 & 0.088 & 0.105\\
SPY 15-min, full & 5{,}075 & 0.171 & 0.016 & 0.037\\
SPY 30-min, full & 2{,}436 & 0.072 & 0.018 & 0.019\\
\addlinespace
WIN 1-min, full & 99{,}812 & 0.201 & 0.056 & 0.056\\
WIN 1-min, first 80\% & 79{,}855 & 0.447 & 0.261 & 0.267\\
WIN 1-min, last 20\% & 19{,}957 & 0.028 & 0.008 & 0.008\\
WIN 5-min, full & 19{,}875 & 0.667 & 0.455 & 0.457\\
WIN 15-min, full & 6{,}565 & 0.850 & 0.983 & 0.981\\
WIN 30-min, full & 3{,}194 & 0.858 & 0.833 & 0.835\\
\bottomrule
\end{tabular}
\end{table}

The short daily SPY sample (2020--2026) confirms the long-sample result. Minute-level evidence is weaker and period-dependent: for both markets the rejection comes mainly from the last 20\% of the sample (September--October 2026), and for WIN it disappears at 5--30-minute sampling. For SPY, aggregation to 15 and 30 minutes strengthens the evidence under the GARCH-type nulls.

\section{Conclusion}\label{sec:conclusion}
An omitted feature does not vanish when an econometric model is fitted to financial returns. Omitted downside jumps move the parameter estimates along a common direction $W$, the score of the omission, which is the same for every model; each model responds by shifting its location, scale and tail parameters to accommodate the extreme downside observations. We showed that this displacement produces a first-order error in value-at-risk forecasts and a second-order regret under the pinball loss for a model that is otherwise correct. In our simulations, omitting jumps on 2\% of the days leaves 99\% VaR 5--7\% too low, leaves regulatory capital 4--5\% below what the true risk requires, and raises the share of Basel yellow-zone outcomes from 11--14\% to 21--22\%.

We proposed a test for one-sided impulses that keeps its nominal size when the fitted model has the wrong tails, scale or volatility dynamics. The local power of the test is governed by the same single-impulse functional that determines the parameter displacement. This link has an uncomfortable implication for risk management: omitted jumps are hardest to detect in exactly the settings where the displacement and the capital shortfall are largest.

Fisher information bounds show that moderate downside jumps fall in an overlap regime, where the asymmetry can be detected but the jump parameters cannot be estimated reliably. We combined these results into a sequential procedure for financial institutions and regulators: test for downside asymmetry, evaluate how much information the sample contains about the jumps, and, in the overlap regime, report the bias of the symmetric model or use a skewed innovation law instead of fitting an unstable jump model. Applied to U.S.\ and Brazilian equity returns, the procedure detects downside asymmetry in all three series, recognizes that it behaves like a skewed innovation law rather than rare jumps, and, for the two U.S.\ series, restores the out-of-sample calibration of 99\% VaR.
% OPEN (D4): optionally one sentence on tail shifts in Tsallis q terms, q=(nu+3)/(nu+1).

\paragraph{Future research: Hausman-type tests for omitted risk.} Corollary~\ref{cor:common} suggests a broader class of tests. Because an omitted feature moves every fitted model along the same direction $W$, but each model responds through its own $\bar H_M^{-1}\bar b_M$, two estimators that agree without the omission separate under it, and their difference estimates $\varepsilon(\bar H_1^{-1}\bar b_1-\bar H_2^{-1}\bar b_2)$. A Hausman-type statistic \citep{hausman1978} based on that difference tests for the omission. The simplest example compares the sample mean with the Student-$t$ location: both estimate the center of a symmetric law even when the tails are wrong, but under downward jumps the mean moves by the full mean jump and the $t$-location much less (Table~\ref{tab:cross}). Such tests need the moments their estimators require and use a single contrast rather than the whole jump direction, so we expect them to be less powerful than the robust test; their value is generality, since the same construction applies to any omitted feature with score $W$. Their theory and power comparisons are left for future research.

\appendix
\section{Three steps behind Theorem~\ref{thm:displacement}}\label{app:steps}
\paragraph{Step 1: the jump-contaminated density is a polynomial in $\varepsilon$.} Conditioning on the set $S\subseteq\{1,\dots,T\}$ of jump dates,
\[
p_\varepsilon(y)=\sum_{S}\varepsilon^{|S|}(1-\varepsilon)^{T-|S|}\,p^{[S]}(y),
\]
where $p^{[S]}$ is the density of $Y$ when jumps occur exactly at the dates in $S$. Each $p^{[S]}$ is a density, so $p_\varepsilon$ is a polynomial in $\varepsilon$ whose coefficients are densities or differences of densities.

\paragraph{Step 2: $\dot p_0=p_0W_T$.} Differentiating at $\varepsilon=0$, only $S=\emptyset$ and $|S|=1$ contribute: $\dot p_0=\sum_t(p^{[t]}-p_0)$. In the HMM, $p^{[t]}$ replaces the state-$j$ emission at date $t$ by $h_j$, so
\[
\frac{p^{[t]}(y)}{p_0(y)}=\E_0\!\left[\frac{h_{K_t}(y_t)}{f_{K_t}(y_t)}\,\Big|\,Y_{1:T}=y\right]=\sum_j\gamma_{tj}(y)\frac{h_j(y_t)}{f_j(y_t)},
\]
and summing over $t$ gives $\dot p_0=p_0W_T$ with $W_T$ as in \eqref{eq:W}.

\paragraph{Step 3: single-impulse representation.} For any integrable statistic $S$, $\E_0[SW_T]=\int S\dot p_0=\sum_t\int S\,(p^{[t]}-p_0)=\sum_t\E_0\E_G\{S(X^{[t,G]})-S(X)\}$, because $p^{[t]}$ is the density of $X^{[t,G]}$.

\section{Proofs}\label{app:proofs}
\paragraph{Proof of Theorem~\ref{thm:largeT}.}

\emph{Step 1: expansion of the expected score.} Let $J_\varepsilon=\{u\le t:B_u=1\}$ be the random set of jump dates. By dominated differentiation, $F(\beta,\varepsilon):=\nabla_\beta q(\beta,\varepsilon)=\E_\varepsilon s_t(\beta)$, and
\[
F(\beta,\varepsilon)-F(\beta,0)=\E\,\Delta_t(J_\varepsilon)
=\sum_{u\le t}\Pr(u\in J_\varepsilon)\,\E\,\Delta_t(\{u\})+R(\beta,\varepsilon),
\]
where we used that $\Delta_t(\{u\})$ depends only on $X$ and $G_u$, which are independent of $B_u$. Hence the first term is $\varepsilon\sum_{j\ge0}\E\Delta_t(\{t-j\})=\varepsilon\,\bar b_M(\beta)$, finite by (iii). By the interaction bound and independence of the $B_u$, $\sup_N\|R(\beta,\varepsilon)\|\le\varepsilon^2\sum_{u<v}c_{t-u}c_{t-v}\le\tfrac12\varepsilon^2(\sum_jc_j)^2$. Thus
\[
F(\beta,\varepsilon)=F(\beta,0)+\varepsilon\,\bar b_M(\beta)+O(\varepsilon^2)\quad\text{uniformly on }N,
\]
and $\bar b_M(\cdot)$ is continuous on $N$ by dominated convergence.

\emph{Step 2: the maximizer moves continuously.} By (iii) applied to $\ell_t$, $\sup_B|q(\beta,\varepsilon)-q(\beta,0)|\le\varepsilon\sum_ja_j$. Since $q(\cdot,0)$ is continuous on the compact $B$ with a unique maximizer, any maximizer of $q(\cdot,\varepsilon)$ converges to $\bar\beta^*_M$ as $\varepsilon\to0$ \citep[Theorem~2.1]{newey1994}. By (iii) applied to $\nabla_\beta s_t$, $\nabla^2_\beta q(\cdot,\varepsilon)$ is within $O(\varepsilon)$ of $\nabla^2_\beta q(\cdot,0)$ uniformly on $N$, so $q(\cdot,\varepsilon)$ is strictly concave on a smaller neighborhood for small $\varepsilon$, and the maximizer there is unique.

\emph{Step 3: linearization.} Write $\Delta\beta=\bar\beta^*_M(\varepsilon)-\bar\beta^*_M$. Since $F(\bar\beta^*_M,0)=0$ and $F(\cdot,0)$ is continuously differentiable with derivative $-\bar H_M$ at $\bar\beta^*_M$, the first-order condition $F(\bar\beta^*_M(\varepsilon),\varepsilon)=0$ and Step~1 give
\[
0=-\bar H_M\Delta\beta+o(\|\Delta\beta\|)+\varepsilon\,\bar b_M(\bar\beta^*_M(\varepsilon))+O(\varepsilon^2).
\]
By Step~2 and continuity of $\bar b_M$, $\bar b_M(\bar\beta^*_M(\varepsilon))=\bar b_M(\bar\beta^*_M)+o(1)$. Hence $\|\Delta\beta\|=O(\varepsilon)$, and then $\Delta\beta=\varepsilon\bar H_M^{-1}\bar b_M+o(\varepsilon)$. The second form of $\bar b_M$ follows from stationarity, by shifting each term $\E\Delta_t(\{t-j\})$ to an impulse at $t$ observed $j$ periods later.

\emph{Step 4: the estimator.} Under the uniform law of large numbers, $\hat\beta_T$ maximises a criterion that converges uniformly to $q(\cdot,\varepsilon)$, whose maximizer is unique for small $\varepsilon$; consistency follows from \citet[Theorem~2.1]{newey1994}.
\hfill$\square$

\paragraph{Proof of Theorem~\ref{thm:local}.}

\emph{Step 1: decomposition at $\beta^*$.} Given $\F^Y_{t-1}$, the current value is $X_t$ with probability $1-\varepsilon_T$ and $X_t-G_t$ with probability $\varepsilon_T$, since $B_t$ is independent of everything else. The past $Y_{1:t-1}$ is a function of $X_{1:t-1}$ and of jump variables independent of $X$. Hence the conditional law of $X_t$ given $\F^Y_{t-1}$ is a mixture, over the conditional law of $X_{1:t-1}$ given $\F^Y_{t-1}$, of the laws of $X_t$ given $X_{1:t-1}$; each is symmetric about the constant $\mu^*$ (Assumption~\ref{A:null}), so the mixture is too. Since $\xi_t(\theta;\beta^*;\cdot;\mathcal P)$ is odd about $\mu^*$ for every past $\mathcal P$,
\[
\E[\xi_t(\theta;\beta^*)\mid\F^Y_{t-1}]=\varepsilon_T\,\psi^{\varepsilon_T}_t(\theta).
\]
Write $\xi_t=m_t+\varepsilon_T\psi^{\varepsilon_T}_t$ with $m_t=\xi_t-\E[\xi_t\mid\F^Y_{t-1}]$, a martingale difference array with respect to $\F^Y_t$. Then
\[
T^{-1/2}\sum_t\xi_t(\theta;\beta^*)=T^{-1/2}\sum_tm_t(\theta)+c\,T^{-1}\sum_t\psi^{\varepsilon_T}_t(\theta).
\]

\emph{Step 2: the drift.} By Assumption~\ref{A:local}(ii), $T^{-1}\sum_t\{\psi^{\varepsilon_T}_t(\theta)-\psi^0_t(\theta)\}\to0$ in $L^1$, uniformly in $\theta$. By ergodicity of $X$, $T^{-1}\sum_t\psi^0_t(\theta)\to\E\psi^0_t(\theta)=\lambda(\theta)$ almost surely, and uniformly in $\theta$ by the Lipschitz condition and compactness of $\Theta$.

\emph{Step 3: the martingale part.} The array $m_t$ is bounded (Assumption~\ref{A:local}(i)), so the conditional Lindeberg condition holds. Its conditional covariance is
\[
\E[m_t(\theta)m_t(\theta')\mid\F^Y_{t-1}]=(1-\varepsilon_T)\,v^{\varepsilon_T}_t(\theta,\theta')+O(\varepsilon_T),
\]
the $O(\varepsilon_T)$ term collecting the jump branch and the squared drift, both bounded. By Assumption~\ref{A:local}(ii) and ergodicity, $T^{-1}\sum_t\E[m_t(\theta)m_t(\theta')\mid\F^Y_{t-1}]\to\E\,v^0_t(\theta,\theta')=\Omega(\theta,\theta')$ in probability. The martingale central limit theorem for arrays \citep[Corollary~3.1]{hall1980} gives the finite-dimensional convergence of $T^{-1/2}\sum_tm_t$ to $\mathbb G$; tightness follows from the bounded Lipschitz constant in $\theta$, as in Theorem~\ref{thm:null}.

\emph{Step 4: estimated parameters.} As in the proof of Theorem~\ref{thm:null},
\[
T^{-1/2}\sum_t\hat\xi_t(\theta)-T^{-1/2}\sum_t\xi_t(\theta;\beta^*)=\Big[T^{-1}\sum_t\nabla_\beta\xi_t(\theta;\bar\beta)\Big]\tr\sqrt T(\hat\beta-\beta^*)-(\hat\kappa-\kappa)\tr T^{-1/2}\sum_ts^\mu_t(\hat\beta).
\]
The conditional mean of $\nabla_\beta\xi_t$ given $\F^Y_{t-1}$ is $(1-\varepsilon_T)$ times its no-jump branch plus $O(\varepsilon_T)$; by Assumption~\ref{A:local}(ii),(iv), ergodicity and Lemma~\ref{lem:mds}, the bracket converges to $\E_0\nabla_\beta\xi_t=0$ uniformly in $\theta$, and with (iii) the first term is $o_p(1)$. The location score $s^\mu_t$ is odd, so Steps 1--3 applied to it give $T^{-1/2}\sum_ts^\mu_t(\beta^*)=O_p(1)$; with $\hat\kappa-\kappa=o_p(1)$ the second term is $o_p(1)$.

\emph{Step 5: studentization.} $T^{-1}\sum_t\hat\xi_t(\theta)^2=\Omega(\theta,\theta)+O_p(\varepsilon_T)+o_p(1)$ by the same arguments, which gives the stated limit of $Z_T$ and, by the continuous mapping theorem, of $\mathcal S_T$.
\hfill$\square$

\paragraph{Proof of Proposition~\ref{prop:jumpinfo}.}

Since $(1-\varepsilon)f+\varepsilon h\ge\varepsilon h$ and $\ge(1-\varepsilon)f$, the integrand in \eqref{eq:Iphi} is at most $(\partial_\varphi h)^2/(\varepsilon h)$ and at most $(\partial_\varphi h)^2/\{(1-\varepsilon)f\}$, which give the two upper bounds; $(1-\varepsilon)f+\varepsilon h\le f+h$ gives the lower bound. For the Student-$t$ case, $|f'(x)|\le c(1+|x|)^{-(\nu+2)}$ and $f(x)\le c(1+|x|)^{-(\nu+1)}$, and $(1+|y|)\le(1+|y+g|)(1+g)$ for $g\ge0$. Hence
\[
|\partial_\theta h(y)|=\big|\E[Zf'(y+\theta Z)]\big|\le c\,\E\big[Z(1+\theta Z)^{\nu+2}\big](1+|y|)^{-(\nu+2)},
\]
\[
|\partial_kh(y)|=\big|\E[f(y+\theta Z)\{\log Z-\psi(k)\}]\big|\le c\,\E\big[(1+\theta Z)^{\nu+1}|\log Z-\psi(k)|\big](1+|y|)^{-(\nu+1)},
\]
with finite expectations because $Z$ has all moments and $\E|\log Z|^2<\infty$. Since $f(y)\ge c'(1+|y|)^{-(\nu+1)}$, the ratios $(\partial_\theta h)^2/f$ and $(\partial_kh)^2/f$ are bounded by multiples of $(1+|y|)^{-(\nu+3)}$ and $(1+|y|)^{-(\nu+1)}$, both integrable. The variance bound is the Cram\'er--Rao bound, using $(I^{-1})_{\varphi\varphi}\ge1/I_{\varphi\varphi}$ when other parameters are estimated jointly.
\hfill$\square$

\section{Additional simulation evidence}\label{app:extra}
\subsection{Size under symmetric non-Student-$t$ contamination: the HMM-9 design}\label{sec:hmm9size}
\begin{table}[ht]\centering\footnotesize\setlength{\tabcolsep}{4pt}\caption{Rejection rates (\%) at the 5\% level in the HMM-9 design, 300 paths per row. Monte Carlo s.e.\ about 1.3 points at 5\%.}\label{tab:hmm9}
\begin{tabular}{llr rrr rrr}
\toprule
 & & & \multicolumn{3}{c}{Static Student-$t$ null} & \multicolumn{3}{c}{GARCH(1,1)-$t$ null}\\
\cmidrule(lr){4-6}\cmidrule(lr){7-9}
Generator & Role & $n$ & Robust & OLS proj. & Naive & Robust & OLS proj. & Naive\\
\midrule
Clean Student-$t$ HMM & size & 1,000 & 4.0 & 4.0 & 0.0 & 4.0 & 4.3 & 0.0\\
 & & 2,000 & 3.7 & 4.0 & 0.3 & 4.0 & 4.0 & 0.0\\
 & & 3,000 & 3.0 & 3.0 & 0.0 & 2.3 & 2.3 & 0.0\\
\addlinespace
Gaussian contamination & size & 1,000 & 3.7 & 3.3 & 0.0 & 3.3 & 3.7 & 0.0\\
 & & 2,000 & 4.3 & 4.3 & 2.0 & 4.7 & 4.7 & 1.7\\
 & & 3,000 & 6.3 & 6.7 & 5.7 & 7.3 & 7.3 & 6.3\\
\addlinespace
Negative Gamma jumps & power & 1,000 & 97.0 & 97.0 & 75.0 & 97.7 & 97.7 & 71.7\\
 & & 2,000 & 100.0 & 100.0 & 99.7 & 100.0 & 100.0 & 99.7\\
 & & 3,000 & 100.0 & 100.0 & 100.0 & 100.0 & 100.0 & 100.0\\
\bottomrule
\end{tabular}
\end{table}
The Gaussian-contamination arm is symmetric but not Student-$t$, so Theorem~\ref{thm:symm} predicts correct size for the robust test. Its rejection rates are 3.3--7.3\%; the highest, 7.3\% at $n=3{,}000$, is about 1.5 Monte Carlo s.e.\ above 5\%. Here the naive score is conservative rather than liberal, and its power at $n=1{,}000$ (72--75\%) is well below that of the robust test (97--98\%). These jumps are large (mean size 6), so this is a high-power design.

\subsection{Numerical checks of the theory}
\paragraph{Orthogonality (Lemma~\ref{lem:mds}).} On a clean 300{,}000-observation path of the HMM-9 design with a static Student-$t$ fit (pseudo-true $\nu^*=3.99$, so the tails are wrong in both regimes), the sample means of $\partial\xi_t/\partial\beta$ for the scale and $\nu$ are within 1.5 standard errors of zero at all five $\theta$, while $\E\,\partial a_t/\partial\mu$ is between 1.09 and 1.98 (standard errors at most 0.024).

\paragraph{Null limit and local power (Theorems~\ref{thm:null}--\ref{thm:local}).} With $T=20{,}000$ from the clean HMM-9 design and 150 paths each: under the null the five $Z_T(\theta)$ have standard deviations 0.98--1.07 and means $-0.17$ to $-0.02$ (s.e.\ 0.08--0.09). Under local jumps with $c=4$ ($\varepsilon_T=0.028$, $G=0.9\,$Gamma$(2,1)$), the first-order prediction of the mean of $Z_T(\theta)$ is $(3.95,5.51,5.82,4.47,2.08)$ and the simulated means are $(3.56,4.92,5.17,4.11,2.23)$. Most of the gap is an $O(\varepsilon)$ studentization effect: under the alternative $T^{-1}\sum\hat\xi_t^2$ estimates $(1-\varepsilon)\Omega+\varepsilon\E[\xi_t^2|_{y_t=X_t-G}]$. With that denominator the prediction is $(3.85,5.16,5.29,4.06,1.98)$, against null-corrected simulated means $(3.73,5.03,5.26,4.19,2.25)$: within 1.2 combined s.e.\ for the first four $\theta$ and 2.2 s.e.\ at $\theta=2.5$.

\paragraph{A practical caveat.} Lemma~\ref{lem:bounded} bounds $a_t$, but the bound grows quickly with $\theta/\sigma_t$: on the long path $\max_t|a_t|$ is 2.4 at $\theta=0.3$ and 1{,}342 at $\theta=2.5$. A grid capped at about $2\hat\sigma$ is a reasonable default.

\section{Reproducibility}
All results can be regenerated from the code and seeds in the public repository \url{https://github.com/profdrgiorno/giorno-tsallis-omitted-risk}, whose README maps every table to the script that produces it. The code is free for research and other noncommercial use. Paths of the HMM-6 manifest are regenerated from \texttt{DGP\_MANIFEST.csv} (repository \texttt{profdrgiorno/\allowbreak Giorno\_garch\_hmm6}); the jump stream is \texttt{default\_rng([seed, 20261031])} and the skewing stream \texttt{default\_rng([seed, 20261103])}. The HMM-9 paths are regenerated from \texttt{simulate()} in \texttt{profdrgiorno/\allowbreak Giorno-Tsallis-Gamma-HMM}. Scripts: \texttt{full\_srj.py} (Table~\ref{tab:full}), \texttt{hmm9\_srj.py} (Table~\ref{tab:hmm9}), \texttt{b3\_var.py} (Table~\ref{tab:var}), \texttt{b12\_run.py} and \texttt{ks\_rerun.py} (Tables~\ref{tab:b12}--\ref{tab:vuong}), \texttt{jumpinfo.py} (Table~\ref{tab:crb}), \texttt{localpower\_check.py}, \texttt{ortho\_check.py}, \texttt{download\_daily\_long.py} and \texttt{procedure.py} (Table~\ref{tab:daily}), \texttt{realdata.py}, \texttt{realdata\_full.py} and \texttt{realdata\_freq.py} (Table~\ref{tab:data}). Raw market data are not redistributed; the daily series are downloaded from Yahoo Finance by \texttt{download\_daily\_long.py}.

% =====================================================================
\begin{thebibliography}{99}
\bibitem[Andrews(1994)]{andrews1994} Andrews, D.W.K. (1994). Empirical process methods in econometrics. In R.F.~Engle and D.L.~McFadden (eds.), \emph{Handbook of Econometrics}, Vol.~4, 2247--2294. Elsevier.
\bibitem[Andrews and Ploberger(1994)]{andrews1994ap} Andrews, D.W.K. and Ploberger, W. (1994). Optimal tests when a nuisance parameter is present only under the alternative. \emph{Econometrica} 62(6), 1383--1414.
\bibitem[Andrews et al.(forthcoming)]{andrews2026} Andrews, I., Barnhard, H. and Carlson, J. (forthcoming). True and pseudo-true parameters. \emph{Econometric Theory}. arXiv:2604.15563.
\bibitem[Andrews et al.(2017)]{andrews2017} Andrews, I., Gentzkow, M. and Shapiro, J.M. (2017). Measuring the sensitivity of parameter estimates to estimation moments. \emph{Quarterly Journal of Economics} 132(4), 1553--1592.
\bibitem[Bai and Ng(2001)]{bainh2001} Bai, J. and Ng, S. (2001). A consistent test for conditional symmetry in time series models. \emph{Journal of Econometrics} 103(1--2), 225--258.
\bibitem[Bai and Ng(2005)]{bainh2005} Bai, J. and Ng, S. (2005). Tests for skewness, kurtosis, and normality for time series data. \emph{Journal of Business \& Economic Statistics} 23(1), 49--60.
\bibitem[Barndorff-Nielsen and Shephard(2006)]{bns2006} Barndorff-Nielsen, O.E. and Shephard, N. (2006). Econometrics of testing for jumps in financial economics using bipower variation. \emph{Journal of Financial Econometrics} 4(1), 1--30.
\bibitem[Basel Committee(1996)]{basel1996} Basel Committee on Banking Supervision (1996). \emph{Supervisory Framework for the Use of ``Backtesting'' in Conjunction with the Internal Models Approach to Market Risk Capital Requirements}. Bank for International Settlements.
\bibitem[Bollerslev(1986)]{bollerslev1986} Bollerslev, T. (1986). Generalized autoregressive conditional heteroskedasticity. \emph{Journal of Econometrics} 31(3), 307--327.
\bibitem[Bonhomme and Weidner(2022)]{bonhomme2022} Bonhomme, S. and Weidner, M. (2022). Minimizing sensitivity to model misspecification. \emph{Quantitative Economics} 13(3).
\bibitem[Chan and Maheu(2002)]{chan2002} Chan, W.H. and Maheu, J.M. (2002). Conditional jump dynamics in stock market returns. \emph{Journal of Business \& Economic Statistics} 20(3), 377--389.
\bibitem[Chen et al.(2024a)]{chen2024cs} Chen, F., Liang, F. and Jiang, Q. (2024a). Conditional symmetry test based on empirical characteristic function. \emph{Journal of Nonparametric Statistics}, doi:10.1080/10485252.2024.2398680.
\bibitem[Chen et al.(2024b)]{chen2024} Chen, X., Hansen, L.P. and Hansen, P.G. (2024b). Robust inference for moment condition models without rational expectations. \emph{Journal of Econometrics} 243, 105653.
\bibitem[Christoffersen and Pelletier(2004)]{christoffersen2004} Christoffersen, P. and Pelletier, D. (2004). Backtesting value-at-risk: a duration-based approach. \emph{Journal of Financial Econometrics} 2(1), 84--108.
\bibitem[Christoffersen(1998)]{christoffersen1998} Christoffersen, P.F. (1998). Evaluating interval forecasts. \emph{International Economic Review} 39(4), 841--862.
\bibitem[Davies(1977)]{davies1977} Davies, R.B. (1977). Hypothesis testing when a nuisance parameter is present only under the alternative. \emph{Biometrika} 64(2), 247--254.
\bibitem[Davies(1987)]{davies1987} Davies, R.B. (1987). Hypothesis testing when a nuisance parameter is present only under the alternative. \emph{Biometrika} 74(1), 33--43.
\bibitem[Ding et al.(1993)]{ding1993} Ding, Z., Granger, C.W.J. and Engle, R.F. (1993). A long memory property of stock market returns and a new model. \emph{Journal of Empirical Finance} 1(1), 83--106.
\bibitem[Douc et al.(2004)]{douc2004} Douc, R., Moulines, E. and Ryd\'en, T. (2004). Asymptotic properties of the maximum likelihood estimator in autoregressive models with Markov regime. \emph{Annals of Statistics} 32(5), 2254--2304.
\bibitem[Duan et al.(2006)]{duan2006} Duan, J.-C., Ritchken, P. and Sun, Z. (2006). Approximating GARCH-jump models, jump-diffusion processes, and option pricing. \emph{Mathematical Finance} 16(1), 21--52.
\bibitem[Eraker et al.(2003)]{eraker2003} Eraker, B., Johannes, M. and Polson, N. (2003). The impact of jumps in volatility and returns. \emph{Journal of Finance} 58(3), 1269--1300.
\bibitem[Escanciano and Olmo(2010)]{escanciano2010} Escanciano, J.C. and Olmo, J. (2010). Backtesting parametric value-at-risk with estimation risk. \emph{Journal of Business \& Economic Statistics} 28(1), 36--51.
\bibitem[Fernandes et al.(2015)]{fernandes2015} Fernandes, M., Mendes, E.F. and Scaillet, O. (2015). Testing for symmetry and conditional symmetry using asymmetric kernels. \emph{Annals of the Institute of Statistical Mathematics} 67(4), 649--671.
\bibitem[Fern\'andez and Steel(1998)]{fernandez1998} Fern\'andez, C. and Steel, M.F.J. (1998). On Bayesian modeling of fat tails and skewness. \emph{Journal of the American Statistical Association} 93(441), 359--371.
\bibitem[Gaglianone et al.(2011)]{gaglianone2011} Gaglianone, W.P., Lima, L.R., Linton, O. and Smith, D.R. (2011). Evaluating value-at-risk models via quantile regression. \emph{Journal of Business \& Economic Statistics} 29(1), 150--160.
\bibitem[Giorno(2026a)]{giornoA} Giorno, J. (2026a). When can market regimes improve value at risk? A proposed information bound for Student-$t$ hidden Markov model identification. Working paper, State University of Maring\'a. doi:10.13140/RG.2.2.12652.53129.
\bibitem[Glosten et al.(1993)]{gjr1993} Glosten, L.R., Jagannathan, R. and Runkle, D.E. (1993). On the relation between the expected value and the volatility of the nominal excess return on stocks. \emph{Journal of Finance} 48(5), 1779--1801.
\bibitem[Gneiting(2011)]{gneiting2011} Gneiting, T. (2011). Making and evaluating point forecasts. \emph{Journal of the American Statistical Association} 106(494), 746--762.
\bibitem[Haas et al.(2004)]{haas2004} Haas, M., Mittnik, S. and Paolella, M.S. (2004). A new approach to Markov-switching GARCH models. \emph{Journal of Financial Econometrics} 2(4), 493--530.
\bibitem[Hall and Heyde(1980)]{hall1980} Hall, P. and Heyde, C.C. (1980). \emph{Martingale Limit Theory and Its Application}. Academic Press.
\bibitem[Hamilton(1989)]{hamilton1989} Hamilton, J.D. (1989). A new approach to the economic analysis of nonstationary time series and the business cycle. \emph{Econometrica} 57(2), 357--384.
\bibitem[Hamilton(1996)]{hamilton1996} Hamilton, J.D. (1996). Specification testing in Markov-switching time-series models. \emph{Journal of Econometrics} 70(1), 127--157.
\bibitem[Hansen(1992)]{hansen1992} Hansen, B.E. (1992). The likelihood ratio test under nonstandard conditions: testing the Markov switching model of GNP. \emph{Journal of Applied Econometrics} 7(S1), S61--S82.
\bibitem[Hansen(1994)]{hansen1994} Hansen, B.E. (1994). Autoregressive conditional density estimation. \emph{International Economic Review} 35(3), 705--730.
\bibitem[Hansen(1996)]{hansen1996} Hansen, B.E. (1996). Inference when a nuisance parameter is not identified under the null hypothesis. \emph{Econometrica} 64(2), 413--430.
\bibitem[Hausman(1978)]{hausman1978} Hausman, J.A. (1978). Specification tests in econometrics. \emph{Econometrica} 46(6), 1251--1271.
\bibitem[Kitamura et al.(2013)]{kitamura2013} Kitamura, Y., Otsu, T. and Evdokimov, K. (2013). Robustness, infinitesimal neighborhoods, and moment restrictions. \emph{Econometrica} 81(3), 1185--1201.
\bibitem[Koenker and Bassett(1978)]{koenker1978} Koenker, R. and Bassett, G. (1978). Regression quantiles. \emph{Econometrica} 46(1), 33--50.
\bibitem[Kou(2002)]{kou2002} Kou, S.G. (2002). A jump-diffusion model for option pricing. \emph{Management Science} 48(8), 1086--1101.
\bibitem[Kuester et al.(2006)]{kuester2006} Kuester, K., Mittnik, S. and Paolella, M.S. (2006). Value-at-risk prediction: a comparison of alternative strategies. \emph{Journal of Financial Econometrics} 4(1), 53--89.
\bibitem[Kupiec(1995)]{kupiec1995} Kupiec, P.H. (1995). Techniques for verifying the accuracy of risk measurement models. \emph{Journal of Derivatives} 3(2), 73--84.
\bibitem[Le Gland and Mevel(2000)]{legland2000} Le Gland, F. and Mevel, L. (2000). Exponential forgetting and geometric ergodicity in hidden Markov models. \emph{Mathematics of Control, Signals, and Systems} 13(1), 63--93.
\bibitem[Lee and Mykland(2008)]{leemykland2008} Lee, S.S. and Mykland, P.A. (2008). Jumps in financial markets: a new nonparametric test and jump dynamics. \emph{Review of Financial Studies} 21(6), 2535--2563.
\bibitem[Lopez(1999)]{lopez1999} Lopez, J.A. (1999). Methods for evaluating value-at-risk estimates. \emph{Federal Reserve Bank of San Francisco Economic Review} 1999(2), 3--17.
\bibitem[Maasoumi and Racine(2009)]{maasoumi2009} Maasoumi, E. and Racine, J.S. (2009). A robust entropy-based test of asymmetry for discrete and continuous processes. \emph{Econometric Reviews} 28(1--3), 246--261.
\bibitem[Maheu and McCurdy(2004)]{maheu2004} Maheu, J.M. and McCurdy, T.H. (2004). News arrival, jump dynamics, and volatility components for individual stock returns. \emph{Journal of Finance} 59(2), 755--793.
\bibitem[Merton(1976)]{merton1976} Merton, R.C. (1976). Option pricing when underlying stock returns are discontinuous. \emph{Journal of Financial Economics} 3(1--2), 125--144.
\bibitem[Newey and McFadden(1994)]{newey1994} Newey, W.K. and McFadden, D. (1994). Large sample estimation and hypothesis testing. In R.F.~Engle and D.L.~McFadden (eds.), \emph{Handbook of Econometrics}, Vol.~4, 2111--2245. Elsevier.
\bibitem[Ramezani and Zeng(2007)]{ramezani2007} Ramezani, C.A. and Zeng, Y. (2007). Maximum likelihood estimation of the double exponential jump-diffusion process. \emph{Annals of Finance} 3(4), 487--507.
\bibitem[Su(2006)]{su2006} Su, L. (2006). A simple test for multivariate conditional symmetry. \emph{Economics Letters} 93(3), 374--378.
\bibitem[Trottier and Ardia(2016)]{trottier2016} Trottier, D.-A. and Ardia, D. (2016). Moments of standardized Fernandez--Steel skewed distributions: applications to the estimation of GARCH-type models. \emph{Finance Research Letters} 18, 311--316.
\bibitem[Vuong(1989)]{vuong1989} Vuong, Q.H. (1989). Likelihood ratio tests for model selection and non-nested hypotheses. \emph{Econometrica} 57(2), 307--333.
\bibitem[White(1982)]{white1982} White, H. (1982). Maximum likelihood estimation of misspecified models. \emph{Econometrica} 50(1), 1--25.
\end{thebibliography}
\end{document}